\documentclass[10pt,a4paper]{article}

% ----------------- Paquetes -------------------%
\usepackage[T1]{fontenc}
\usepackage[utf8]{inputenc}
\usepackage[a4paper,margin=2.5cm]{geometry}
\usepackage{mathptmx}             
\usepackage{changepage} 
\usepackage{setspace}
\usepackage[spanish]{babel}
\usepackage[labelfont=bf,labelsep=space]{caption}
\usepackage{amssymb}
\usepackage{amsmath}
\usepackage{ebgaramond}
\usepackage{placeins}
\usepackage{etoolbox}
\usepackage{setspace}
\usepackage{parskip} 
\usepackage{tcolorbox}
\usepackage{needspace}
\usepackage{xparse}
\usepackage{ocgx2}
\usepackage{hyperref}
\usepackage{hypcap}
\usepackage{soul}\sethlcolor{yellow}% resaltado amarillo: \hl{texto} (Panel TCEE, ⌃H)



%%%%%%%%%%%%%%%%%%%%%%%%%%%%%%%%%%%%%%%%%%%%%%%%%%%%%%%%%%%%%%%%
% --- Fija primero tus valores globales "buenos" ---
\setstretch{1.2}                 % Interlineado global
\setlength{\parskip}{0.7\baselineskip}
\setlength{\parindent}{0pt}
\newlength{\GlobalParskip}
\setlength{\GlobalParskip}{\parskip}

\newcounter{eqnum}[subsection]
\renewcommand{\theeqnum}{\arabic{eqnum}}
\newcounter{alphaitem}
\newcounter{numitem}

%% ------------------- Recuadros----------------------%%
\newcommand{\puntosclave}[1]{%
  \begin{tcolorbox}[
    colframe=black,
    title={Puntos clave},
    before upper={\setlength{\parindent}{0.5cm}\setlength{\parskip}{0.5\baselineskip}}
  ]
  #1
  \end{tcolorbox}
}

\newcommand{\modificaciones}[1]{
  \begin{tcolorbox}[
    colback=white,
    colframe=red,
    title={Modificaciones a realizar},
    before upper={\setlength{\parindent}{0.5cm}\setlength{\parskip}{0.5\baselineskip}}
  ]
  #1
  \end{tcolorbox}
}

%% ------------------- Preguntas TEST----------------------%%
% Opciones: radio (single-choice)
\NewDocumentCommand{\OptRadio}{m m m}{%
  \noindent\CheckBox[radio,name=#1,value=#2,width=1.2ex,height=1.2ex]{}%
  \;\textbf{#2.}~#3\par
}

% Opciones: checkbox (multi-choice)
\NewDocumentCommand{\OptCheck}{m m}{%
  \noindent\CheckBox[name=#1,width=1.2ex,height=1.2ex,bordercolor={0 0 0}]{}%
  \;#2\par
}

% Pregunta single-choice (radio)
% Args: 1=id, 2=enunciado, 3=opciones, 4=solución+justificación, 5=imagen (o {}), 6=origen
\NewDocumentCommand{\PreguntaTestSingle}{m m m m m m}{%
  \par\medskip
  \noindent\textbf{#1.}~#2\par\smallskip

    % Imagen (si no es {})
    \IfBlankTF{#5}{}{%
      \begin{center}
        \includegraphics[width=0.75\linewidth]{#5}
      \end{center}
    }

    \begin{Form}
      #3
    \end{Form}

    \par\smallskip
    \switchocg{sol-#1}{\fbox{Mostrar / ocultar solución}}%
    \hfill{\footnotesize #6}\par\smallskip

    \begin{ocg}{Solución}{sol-#1}{0}
      \noindent #4
    \end{ocg}
  \par\medskip
}

% Pregunta multi-choice (checkboxes)
\NewDocumentCommand{\PreguntaTestMulti}{m m m m m m}{%
  \par\medskip
  \noindent\textbf{#1.}~#2\par\smallskip

    % Imagen (si no es {})
    \IfBlankTF{#5}{}{%
      \begin{center}
        \includegraphics[width=0.75\linewidth]{#5}
      \end{center}
    }
    \begin{Form}
      #3
    \end{Form}

    \par\smallskip
    \switchocg{sol-#1}{\fbox{Mostrar / ocultar solución}}%
    \hfill{\footnotesize #6}\par\smallskip

    \begin{ocg}{Solución}{sol-#1}{0}
      \noindent #4
    \end{ocg}
  \par\medskip
}

%% ------------------- CITA ----------------------%%
% \begin{cita}[Autor][Año][Obra] texto \end{cita}: cita sangrada y, debajo a la derecha, «AUTOR (Año), Obra». Los tres datos son opcionales.
% En el Panel TCEE: escribir \cita e Intro (el cursor va al texto; Tab pasa a Autor, Año y Obra).
\NewDocumentEnvironment{cita}{ O{} O{} O{} +b }{%
  \begin{adjustwidth}{2cm}{0pt}
  \raggedright
  #4\par
  \raggedleft
  \if\relax\detokenize{#1#2#3}\relax
    % sin firma
  \else
    #1%
    \if\relax\detokenize{#2}\relax\else\if\relax\detokenize{#1}\relax\else\space\fi(#2)\fi
    \if\relax\detokenize{#3}\relax\else\if\relax\detokenize{#1#2}\relax\else, \fi\textit{#3}\fi
  \fi
  \end{adjustwidth}
}{}



%% ------------------- Listas----------------------%%
    % --- Lista alfabética ---
    \newenvironment{la}
      {\begin{list}{\alph{alphaitem})}{%
          \usecounter{alphaitem}%
          \setlength{\leftmargin}{1cm}%
          \setlength{\parskip}{3pt}% 
          \setlength{\itemsep}{0pt}% ← espacio entre ítems
          \setlength{\parsep}{0pt}%  ← párrafos dentro de un ítem
      }}
      {\end{list}}
      
    % --- Lista numérica ---
    \newenvironment{lnum}
      {\begin{list}{\arabic{numitem}.}{%
          \usecounter{numitem}%
          \setlength{\leftmargin}{1cm}%
          \setlength{\parskip}{3pt}% 
          \setlength{\itemsep}{0pt}% ← espacio entre ítems
          \setlength{\parsep}{0pt}%  ← párrafos dentro de un ítem
      }}
      {\end{list}}

%% ------------------- Preguntas Test ----------------------%%
      \usepackage{xparse} % si ya lo tienes, no lo repitas

    \NewDocumentEnvironment{test}{m m o}{%
      \begin{tcolorbox}[
        colback=white,
        colframe=black!15,
        boxrule=0.6pt,
        arc=2mm,
        left=2mm,right=2mm,top=2mm,bottom=2mm
      ]
      \centering
      \includegraphics[width=\linewidth]{#1}
      \end{tcolorbox}
    
      \vspace{2mm}
    
      % Caja SOLO para texto (SÍ partible y mantiene márgenes al partir)
      \begin{tcolorbox}[
        enhanced,
        breakable,
        colback=black!3,
        colframe=black!25,
        boxrule=0.6pt,
        arc=2mm,
        left=2mm,right=2mm,top=1.5mm,bottom=1.5mm
      ]
      \textbf{Respuesta:} \textbf{#2}%
      \IfValueT{#3}{\par\smallskip \textbf{Justificación:} #3}
      \end{tcolorbox}
    }{}
      
% ------------------- Imágenes----------------------%
    \graphicspath{{Img/}}% Rutas por defecto 
    % --- Imagen adaptativa ---
    % Registros de dimensión definidos una sola vez
    \newdimen\imgcap
    \newdimen\imgmin
    \newdimen\imgmax
    \newdimen\imgremain
    \newdimen\imgh
    \newdimen\imgneed
    
    \newcommand{\imagenfit}[3]{%
      \addvspace{10pt}% separación antes
      \par\begingroup
        % Parámetros de composición (asignaciones, no \newdimen)
        \imgcap=1\baselineskip            % alto estimado de título+fuente
        \imgmin=0.15\textheight
        \imgmax=0.15\textheight
        \imgremain=\pagegoal
        \ifdim\pagegoal=\maxdimen \imgremain=0pt\fi
        \advance\imgremain by -\pagetotal
        \advance\imgremain by -\imgcap
        % Decisión de altura destino
        \ifdim\imgremain<\imgmin
          \newpage
          \imgh=0.20\textheight
        \else
          \imgh=\imgremain
          \ifdim\imgh>\imgmax \imgh=\imgmax \fi
        \fi
        % Reservar espacio para evitar cortes del bloque completo
        \imgneed=\imgh \advance\imgneed by \imgcap
        \Needspace{\imgneed}
        % Cuerpo
        \noindent\begin{minipage}{\textwidth}
          \centering
          \captionof{figure}{\textemdash\ #2}\par\vspace{0.3em}% título arriba
          \includegraphics[width=\linewidth,height=\imgh,keepaspectratio]{\detokenize{#1}}\par
          \vspace{0.3em}\small\textit{Fuente: }#3% fuente abajo
        \end{minipage}%
      \endgroup
      \par\addvspace{10pt}%
    }

%------------ Bloque de RECUADRO ESPECIAL -------------%
    \newenvironment{specialblock}[1]{%
      \begin{quote} % crea sangría a izquierda y derecha
      \small % texto más pequeño
      \noindent\textbf{#1}\\[-0.3em] % título en negrita
      \rule{\linewidth}{0.4pt}\\[0.6em]}{
      \end{quote}}
      
%-------------------------- Portada -----------------------%
\newcommand{\oldtitle}[1]{\def\OldTitle{#1}}
\newcommand{\changerationale}[1]{\def\ChangeWhy{#1}}

\newcommand{\changebox}{%
\begin{tcolorbox}[title={Historial de cambios del tema}]
\textbf{Título anterior:} \OldTitle\par
\textbf{Motivo del cambio:} \ChangeWhy
\end{tcolorbox}
}

%-------------------------- Author Font -----------------------%
    \newcommand{\authorfont}[1]{{\fontfamily{EBGaramond-TLF}\selectfont #1}}

%-------------------------- Anexos -----------------------%
    \makeatletter
    \newcounter{anexo}
    \renewcommand{\theanexo}{\Roman{anexo}}
    \newcommand{\anexo}[1]{%
      \clearpage
      \phantomsection
      \refstepcounter{anexo}%
      \noindent\textbf{Anexo \theanexo.- }#1\par\medskip
      \addcontentsline{stc}{section}{Anexo \theanexo.- #1}%
    }
    \makeatother

    %-------------------------- Lista de figuras (personal) -----------------------%
\makeatletter
\newcommand{\MyListOfFigures}{%
  \clearpage
  \phantomsection
  {\centering\bfseries\Large Índice de figuras\par}%
  \medskip
  % Entrada en nuestro índice de contenido (stc)
  \addcontentsline{stc}{section}{Índice de figuras}%
  % Recuperación del listado (sin cabecera propia)
  \begingroup
    \parindent=0pt\parskip=2pt
    \@starttoc{lof}%
  \endgroup
}
\makeatother

%-------------------------- Bibliografía -----------------------%
\makeatletter
% Contador para las referencias
\newcounter{myref}
% Cabecera de la bibliografía, entra en el índice 'stc'
\newcommand{\MyBibliography}{%
  \clearpage
  \phantomsection
  {\centering\bfseries\Large Bibliografía y referencias\par}%
  \medskip
  \addcontentsline{stc}{section}{Bibliografía y referencias}%
  \setcounter{myref}{0}%
}
\newcommand{\MyBibItem}[4]{%
  \refstepcounter{myref}%
  \noindent[\themyref]\ #1.\ \emph{#2}. #3. #4\par
}
\makeatother

%--------- Establecer cuatro niveles de índice --------%
    \setcounter{tocdepth}{4}   
    \setcounter{secnumdepth}{4}% numera hasta \paragraph

%%%%%%%%%%%%%%%%%%%%%%%%%%%%%%%%

\makeatletter

    %--------- Interlineado Global --------%
    \edef\GlobalBaseLineStretch{\baselinestretch}
    
    %--------- Espaciado de seccinones --------%
    \renewcommand\section{\@startsection{section}
      {1}{\z@}
      {2.5ex \@plus 1ex \@minus .2ex} % espacio antes
      {1.5ex \@plus .2ex}             % espacio después
      {\normalfont\Large\bfseries}    % formato del título
    }
    \renewcommand\subsection{\@startsection{subsection}{2}{\z@}
      {3ex \@plus 0.8ex \@minus .2ex}
      {1.5ex \@plus .2ex}
      {\normalfont\large\bfseries}
    }
    \renewcommand\subsubsection{\@startsection{subsubsection}{3}{\z@}
      {3ex \@plus 0.5ex \@minus .2ex}
      {1ex \@plus .2ex}
      {\normalfont\normalsize\bfseries}
    }
    \renewcommand\paragraph{\@startsection{paragraph}{4}{\z@}%
    {-2ex \@plus -0.5ex \@minus -.2ex}% espacio antes
    {0.8ex \@plus .2ex}% espacio después
    {\normalfont\normalsize\itshape}}% ← cursiva en lugar de negrita

    %--------- Ecuaciones --------%   
    % Variables globales para almacenar títulos y notas
    \gdef\leq@current@section{}%
    \gdef\leq@current@subsection{}%
    \gdef\leq@last@printed@section{}%
    \gdef\leq@last@printed@subsection{}%
    
    % Comando para añadir nota (invisible en el cuerpo)
    \newcommand{\eqnote}[1]{%
      \addcontentsline{leq}{leqnote}{#1}%
    }
    
    % Comando para ecuaciones
    \newcommand{\eqblock}[2]{%
      % Verificar si necesitamos imprimir el título de sección
      \ifx\leq@current@section\leq@last@printed@section\else
        \addcontentsline{leq}{leqsection}{\leq@current@section}%
        \global\let\leq@last@printed@section\leq@current@section
        \gdef\leq@last@printed@subsection{}% Reset subsección al cambiar sección
      \fi
      % Verificar si necesitamos imprimir el título de subsección
      \ifx\leq@current@subsection\leq@last@printed@subsection\else
        \ifx\leq@current@subsection\@empty\else
          \addcontentsline{leq}{leqsubsection}{\leq@current@subsection}%
          \global\let\leq@last@printed@subsection\leq@current@subsection
        \fi
      \fi
      % Ecuación
      \refstepcounter{eqnum}%
      \begin{equation*}
        #1\tag{E.\theeqnum}
      \end{equation*}%
      \addcontentsline{leq}{leqt}{[E.\theeqnum] -- #2}%
      \addcontentsline{leq}{leqm}{#1}%
    }
    
    %%% ----- Índice de ecuaciones ----------%%%
    % Tamaño para []
    \newcommand{\leqIndexFont}{\normalsize}
    \newcommand{\leqTitleFont}{\tiny}
    \newcommand{\leqNoteFont}{\tiny}
    % Cabecera de sección 
    \newcommand*\l@leqsection[2]{%
      \addvspace{25pt}% Mayor separación antes de nueva sección
      \noindent\leqIndexFont
      \makebox[\linewidth]{\rule{.38\linewidth}{0.4pt}\ \ \textbf{  \MakeUppercase{#1}}\ \ \rule{.38\linewidth}{0.4pt}}%
      \par\vspace{-10pt}
    }
    % Cabecera de subsección 
    \newcommand*\l@leqsubsection[2]{%
      \par\addvspace{15pt}%
      \noindent\leqIndexFont #1
      \par\vspace{-7pt}
    }
    % Título de ecuación 
    \newcommand*\l@leqt[2]{%
    \par\addvspace{10pt}%
      \par\noindent\leqTitleFont #1\par
    }
    % Miniatura de ecuación
    \newcommand*\l@leqm[2]{%
      \vspace{0.3\baselineskip}%
      \par\scriptsize\noindent\hspace*{1.5em}\(#1\)\par%
    }
    % Nota de ecuación 
    \newcommand*\l@leqnote[2]{%
      \vspace{0.2\baselineskip}%
      \par\noindent\leqNoteFont\hspace*{1.5em}\textbf{Nota.--} #1\par\vspace{0.5\baselineskip}%
    }
    % Generación del índice
    \newcommand{\mylistofequations}{%
      \clearpage
      \phantomsection
      % Título visual de la lista (ajústalo a tu gusto):
      {\centering\bfseries\Large Índice de ecuaciones\par}
      % Entrada en nuestro índice 'stc':
      \addcontentsline{stc}{section}{Índice de ecuaciones}%
      % Llamada a tu lista real (usa el nombre exacto de tu macro):
      \@starttoc{leq}%
    }
    
    %--------- Captura de títulos de sección --------%
    \let\leq@original@section\section
    \let\leq@original@subsection\subsection
    % Flag para saber si estamos en una sección asterisco
    \newif\ifLeqStarSection
    \LeqStarSectionfalse
    
    % Redefinir \section CON CONTROL DE ASTERISCO
    \renewcommand{\section}{%
      \@ifstar{\leq@section@star}{\@dblarg\leq@section@nostar}%
    }
    \def\leq@section@star#1{%
      \global\LeqStarSectiontrue
      \gdef\leq@current@section{#1}%
      \gdef\leq@current@subsection{}%
      \leq@original@section*{#1}%
      \global\LeqStarSectionfalse
    }
    \def\leq@section@nostar[#1]#2{%
      \global\LeqStarSectionfalse
      \gdef\leq@current@section{#2}%
      \gdef\leq@current@subsection{}%
      \leq@original@section[#1]{#2}%
    }
    % Redefinir \subsection
    \renewcommand{\subsection}{%
      \@ifstar{\leq@subsection@star}{\@dblarg\leq@subsection@nostar}%
    }
    \def\leq@subsection@star#1{%
      \global\LeqStarSectiontrue
      \gdef\leq@current@subsection{#1}%
      \leq@original@subsection*{#1}%
      \global\LeqStarSectionfalse
    }
    \def\leq@subsection@nostar[#1]#2{%
      \global\LeqStarSectionfalse
      \gdef\leq@current@subsection{#2}%
      \leq@original@subsection[#1]{#2}%
    }

    % ----- Restauración de valores Globales de estilo ----------%
    % Flags
    \newif\ifInFootnote
    \InFootnotefalse
    \pretocmd{\@footnotetext}{\InFootnotetrue}{}{}
    \apptocmd{\@footnotetext}{\InFootnotefalse}{}{}
    % Profundidad de listas (soporta anidadas)
    \newcount\MyListDepth \MyListDepth=0
    \AtBeginEnvironment{la}{\advance\MyListDepth by 1}
    \AfterEndEnvironment{la}{\advance\MyListDepth by -1}
    \AtBeginEnvironment{lnum}{\advance\MyListDepth by 1}
    \AfterEndEnvironment{lnum}{\advance\MyListDepth by -1}
    \AtBeginEnvironment{itemize}{\advance\MyListDepth by 1}
    \AfterEndEnvironment{itemize}{\advance\MyListDepth by -1}
    \AtBeginEnvironment{enumerate}{\advance\MyListDepth by 1}
    \AfterEndEnvironment{enumerate}{\advance\MyListDepth by -1}
    % Restauración diferida: hazlo al INICIO del próximo párrafo real
    \newif\if@pendingRestore
    \@pendingRestorefalse
    \newcommand{\RequestRestore}{\global\@pendingRestoretrue}
    \everypar{%
      \if@pendingRestore
        \global\@pendingRestorefalse
        \ifInFootnote\else
          \ifnum\MyListDepth=0
            \setlength{\parskip}{\GlobalParskip}%
            \setstretch{\GlobalBaseLineStretch}%
          \fi
        \fi
      \fi
    }
    % Al final de bloques:
    \AfterEndEnvironment{equation}{\RequestRestore}
    \AfterEndEnvironment{equation*}{\RequestRestore}
    \AfterEndEnvironment{la}{\RequestRestore}
    \AfterEndEnvironment{lnum}{\RequestRestore}
    % Al final de macros:
    \apptocmd{\eqblock}{\RequestRestore}{}{}
    \apptocmd{\imagenfit}{\RequestRestore}{}{}
    
\makeatother

% ===================== ToC propio con 4 niveles (único bloque) =====================
\makeatletter

% --- Escritores: añadimos una línea al .stc en cada cabecera numerada ---
% Guardamos las versiones actuales para llamar al comportamiento normal
\let\MyTOC@orig@section\section
\let\MyTOC@orig@subsection\subsection
\let\MyTOC@orig@subsubsection\subsubsection
\let\MyTOC@orig@paragraph\paragraph

% SECTION
\renewcommand{\section}{%
  \@ifstar{\MyTOC@section@star}{\@dblarg\MyTOC@section@nostar}%
}
\def\MyTOC@section@star#1{%
  \MyTOC@orig@section*{#1}% sin entrada al .stc
}
\def\MyTOC@section@nostar[#1]#2{%
  \MyTOC@orig@section[{#1}]{#2}% comportamiento normal (escribe al .toc)
  % Escribimos también nuestra línea al .stc (texto con \numberline)
  \addcontentsline{stc}{section}{\protect\numberline{\thesection}#2}%
}

% SUBSECTION
\renewcommand{\subsection}{%
  \@ifstar{\MyTOC@subsection@star}{\@dblarg\MyTOC@subsection@nostar}%
}
\def\MyTOC@subsection@star#1{%
  \MyTOC@orig@subsection*{#1}%
}
\def\MyTOC@subsection@nostar[#1]#2{%
  \MyTOC@orig@subsection[{#1}]{#2}%
  \addcontentsline{stc}{subsection}{\protect\numberline{\thesubsection}#2}%
}

% SUBSUBSECTION
\renewcommand{\subsubsection}{%
  \@ifstar{\MyTOC@subsubsection@star}{\@dblarg\MyTOC@subsubsection@nostar}%
}
\def\MyTOC@subsubsection@star#1{%
  \MyTOC@orig@subsubsection*{#1}%
}
\def\MyTOC@subsubsection@nostar[#1]#2{%
  \MyTOC@orig@subsubsection[{#1}]{#2}%
  \addcontentsline{stc}{subsubsection}{\protect\numberline{\thesubsubsection}#2}%
}

% PARAGRAPH
\renewcommand{\paragraph}{%
  \@ifstar{\MyTOC@paragraph@star}{\@dblarg\MyTOC@paragraph@nostar}%
}
\def\MyTOC@paragraph@star#1{%
  \MyTOC@orig@paragraph*{#1}%
}
\def\MyTOC@paragraph@nostar[#1]#2{%
  \MyTOC@orig@paragraph[{#1}]{#2}%
  % Si no numeras los \paragraph, cambia \theparagraph por texto plano sin \numberline
  \addcontentsline{stc}{paragraph}{\protect\numberline{\theparagraph}#2}%
}

% --- Impresor del índice propio (lee el .stc) ---
\newcommand{\MyTOC}{%
  \par\bigskip
  {\centering\bfseries\Large Índice de contenido\par}%
  \medskip
  \begingroup
    \parindent=0pt\parskip=2pt
    \@starttoc{stc}%
  \endgroup
}

\makeatother
% ===================== fin ToC propio =====================


%%%%%%%%%%%%%%


% --- Datos del tema ---
\title{Tema 3.A.17 -- Análisis de mercados (II)\\
\large Teoría del monopolio. Discriminación de precios. Monopolio natural. Producción conjunta. Análisis de eficiencia y bienestar. Monopsonio. Monopolio bilateral.}
\author{Marcelo Ortega}
\date{Febrero 2026}
\newcommand{\TemaCode}{3A17}

\oldtitle{Tema 3.A.15. Teoría del monopolio. Regulación y control.}
\changerationale{
    El tema prescinde de la parte de regulación y control, que pasa a tratarse en el nuevo tema 20 bajo una óptica más exhaustiva y coherente; la regulación en presencia de poder de mercado. En este tema el espacio que deja el contenido eliminado debería de suplirse con una mayor profundización en los desarrollos teóricos conectados al monopolio. De ahí que el título haga explícitos algunos de esos contenidos en los que se espera un tratamiento adecuado por parte del opositor.
}

\begin{document}


\maketitle
\changebox

\newpage

% --- Página de resumen y modificaciones ---
\puntosclave{ %% Escribir puntos clave Apuntes CECO
    }
\vspace{1cm}

\modificaciones{
    
    }

\newpage
\MyTOC
\newpage

\mylistofequations
\clearpage

\MyListOfFigures
\clearpage


\pagenumbering{arabic}
\setcounter{page}{1}


\section{Introducción}\label{intro}


\subsection{Contextualización}\label{context}


\subsubsection{Relevancia}\label{importance}


\subsubsection{Historia}\label{history}



\subsubsection{Problemática}\label{problem} 




\subsection{Estructura de la exposición}\label{structure}






\clearpage
\section{Teoría del Monopolio}
\puntosclave{
    En el monopolio uniproducto, el equilibrio de mercado viene dado por la maximización del beneficio por parte del monopolista, siéndole indiferente para ello elegir precio o cantidad. Esta maximización se produce cuando igualan sus ingresos y costes marginales.
    
	No existe una única curva de oferta del monopolista que pueda derivarse de su curva de costes marginales.
    
	El precio es mayor (y el output menor) en monopolio (uniproducto) que en competencia perfecta, produciéndose ineficiencia en la asignación. La pérdida de eficiencia que se produce bajo el monopolio recibe el nombre de peso muerto (deadweight loss).
    
	En presencia de restricciones de capacidad, el monopolista (uniproducto) debe tener el número suficiente de plantas para producir la cantidad óptima. La decisión de asignación de dicha producción dependerá de si los costes marginales de producción en cada planta son o no constantes.
}

La teoría microeconómica distingue distintas estructuras de mercado según el grado de interdependencia estratégica entre las empresas presentes por el lado de la oferta en el mercado de referencia, lo cual determina la competencia entre ellas: competencia perfecta, monopolio, oligopolio y competencia monopolística.\footnote{%
    Por interdependencia estratégica cabe entender la noción de que cada empresa es consciente de que sus beneficios no sólo dependen de sus decisiones sobre las variables competitivas, sino también de las decisiones tomadas a este respecto por las demás empresas presentes en el mercado. Las situaciones económicas con multiplicidad de agentes varían enormemente en su grado de interdependencia estratégica:
\begin{la}
    \item En situaciones de monopolio (donde solo una empresa ofrece el bien objeto de análisis) o de competencia perfecta (donde todos los agentes actúan como precio aceptantes), la naturaleza de la interdependencia estratégica es prácticamente nula (por ello, no es necesario el empleo de la teoría de juegos); y,
    \item Sin embargo, en otras situaciones como los mercados oligopólicos (donde hay más de un vendedor de un bien, pero no multitud de ellos), la interdependencia estratégica fue juega un papel central (que hace indispensable el empleo de la teoría de juegos para su análisis)
\end{la}
}

El objeto de nuestra exposición consistirá en estudiar la determinación del precio y la cantidad producida en equilibrio en una estructura de mercado monopolística, de acuerdo con los resultados que ofrece la teoría de la organización industrial.\footnote{%
    Mientras que en Estados Unidos se habla de organización industrial, en Europa se halla más extendida la denominación de economía industrial.
}

Un mercado monopolístico es la estrucura de mercado que presenta la característica fundamental de tener un único productor de un bien o conjunto de bienes que no cuenta con sustitutos próximos. Esta situación viene dada por la existencia de barreras de entrada cuyo origen puede estar en:
\begin{la}
    \item Comportamientos estratégicos por parte del monopolista que impiden la entrada de potenciales competidores; o,
    \item Restricciones legales o administrativas que conceden una situación de privilegio monopolista a una determinada empresa (ejemplos: la concesión gubernamental de licencias o la imposición de barreras comerciales que excluyan a competidores extranjeros, derechos de patente sobre un producto o proceso productivo, etc.); o,
    \item Una tecnología de producción que propicie la subaditividad de costes de un único productor, determinando así una situación de monopolio natural.\footnote{%
        Existen otros motivos residuales que pueden dar lugar a la existencia del monopolio, como:
    \begin{lnum}
        \item La propiedad de materias primas estratégicas; o,
        \item El tamaño reducido del mercado (p.e. la tienda/bar de un pueblo)
    \end{lnum}
    }
\end{la}
	
Como consecuencia de estas características, la empresa monopolista tiene poder de mercado y por tanto, da lugar a un \textbf{\textit{fallo de mercado}}.\footnote{%
    Al ser única la empresa que vende el producto, y suponiendo que la función de demanda de mercado es continua y decreciente en el precio:
\begin{lnum}
    \item La cantidad de output que la empresa venda responderá de forma continua ante variaciones del precio que fije la empresa;
    \item De modo equivalente, el precio que fije la empresa responderá de forma continua ante variaciones del output que venda.
\end{lnum}
El monopolio ya era rentable aumentar su precio por encima del nivel competitivo porque no conduce a una pérdida completa de sus ventas.
}

Para estudiar el equilibrio de mercado bajo monopolio, se empleará un contexto especial pero analíticamente muy tratable: el modelo de equilibrio parcial marshalliano;\footnote{%
    el tamaño reducido del mercado permite dos simplificaciones relevantes de cara al análisis del equilibrio en el mercado:
\begin{lnum}
    \item Cuando el gasto en el bien bajo consideración es una proporción pequeña del gasto total del consumidor, solo una fracción reducida de una unidad monetaria adicional de renta se gastará en ese bien. Por lo tanto, cabe esperar que los efectos renta para este bien sean reducidos; y,
    \item De modo similar, con efecto sustitución dispersos, el pequeño tamaño del mercado objeto de estudio debería conducir a que los precios de los otros bienes se vieran casi inalterados ante cambios en el mercado considerado. Debido a las fijeza de los otros precios, queda justificado considerar el gasto en nuestros otros bienes como un único bien compuesto, llamado Numerario.
\end{lnum}
} e iniciaremos el análisis considerando el caso más sencillo: ¿Cómo un monopolio uniproducto determina el precio y la cantidad que maximiza su beneficio?

\subsection{Monopolio uniproducto}
Supongamos que ambos agentes maximizan su utilidad (con los ajustes propios de la variable que se lo proporciona, i.e. consumidores, cantidad; productores, beneficios), que estamos en un entorno de información completa acerca de las características tanto del producto como de la demanda de éste y que la función de demanda agregada es derivable y decreciente en el precio (por tanto, la función inversa de demanda tiene pendiente negativa: $P^d_x(x)$).\footnote{%
    Por lo tanto, se trata de un bien normal. Además, Cabe señalar que siguiendo nuestro enfoque de equilibrio parcial marshalliano, se han supuesto funciones de utilidad cuasi lineales de los consumidores para alcanzar esta función de demanda en la que no hay efecto renta ni efecto sustitución cruzados.
} 

Para obtener el equilibrio de mercado, basta con resolver un problema de optimización consistente en maximizar el beneficio, bien sea en función del precio (p: fijación del precio que maximiza beneficios), como en función de la cantidad a producir (q: fijación de la cantidad ofertada que maximice beneficios).\footnote{%
    Bajo monopolio resulta equivalente elegir precio o cantidad, dado que precio y el output están relacionados de forma unívoca a través de la función de demanda. Aunque se suponga que el monopolista fija el precio y los consumidores eligen la cantidad como función del precio, cabe concebir que el monopolista elige la cantidad óptima que quiere que los consumidores adquieran y, a continuación, fija el precio correspondiente
}

\subsubsection{Desarrollo formal}
Si consideramos el ajuste en cantidades, el problema queda como sigue:

\eqblock{
\begin{aligned}
    &\boxed{
    \begin{aligned}
        &\arg\max_{Q} &&\Pi_M(Q)=\underbrace{\underbrace{P_x^d(Q)}_{\substack{\text{\scriptsize Demanda residual}\\\text{\scriptsize es demanda de mercado}}}\cdot Q}_{IT(Q)\equiv \text{\scriptsize Ingresos Totales}}-C(Q)\\[0.5em]
        &\text{s.a.} &&\Pi_M(Q)\geq\Pi_M(0)
    \end{aligned}
    }\\[1em]
    &\text{CPO (Condiciones necesarias):}\\[0.5em]
    &\underbrace{\frac{\partial}{\partial Q}\Pi_M(Q)=0\Longrightarrow\frac{\partial}{\partial Q}\Pi_M(Q)\quad=\quad\frac{\partial}{\partial Q}P_x^s(Q)\cdot Q+P_x^s(Q)-\frac{\partial}{\partial Q}C(Q)=0}_{\text{Eq:}\;(Q^*)\;\boxed{{IMg}_x(Q)={CMg}_x(Q)}}\\[1em]
    &\text{CSO (Condiciones suficientes)}:\\[0.5em]
    &\frac{\partial}{\partial Q}\Pi_M(Q)<0\;\Rightarrow\;\frac{\partial}{\partial Q}{IMg}_x(Q)-\frac{\partial}{\partial Q}{CMg}_x(Q)<0\;\Rightarrow\;\boxed{\frac{\partial}{\partial Q}{IMg}_x(Q)<\frac{\partial}{\partial Q}{CMg}_x(Q)}
\end{aligned}
}{Problema de maximización: Monopolio uniproducto. Teoría del Monopolio}

Por tanto, si el monopolio maximiza sus beneficios, estará eligiendo directamente la cantidad que desea producir sabiendo ($\sim$ internalizando) que el precio al que venderá esa cantidad viene determinado por la curva de demanda del mercado. A diferencia de la empresa competitiva, el monopolista no toma el precio como dado, sino que internaliza que cualquier variación en la cantidad ofrecida altera el precio de mercado. Por ello, no existe una curva de oferta que relacione de forma independiente precio y cantidad a partir de los costes: la decisión óptima del monopolista surge exclusivamente del problema de maximización, donde la condición relevante es la igualdad entre Ingreso Marginal y Coste Marginal.

En este contexto, el Coste Marginal ya no determina el precio, sino únicamente la cantidad óptima a producir. El precio se fija posteriormente, como el máximo que los consumidores están dispuestos a pagar por esa cantidad, esto es, evaluando dicha cantidad en la curva de demanda. De este modo, la oferta del monopolista no es una función tecnológica, sino una elección estratégica condicionada por la demanda.  

Analizando los componentes del problema de maximización podemos encontrar:
\begin{la}
    \item \textbf{Función Inversa de Demanda Agregada} ($P_x^d(x)$). La curva de demanda conocida por el monopolio dará una relación unívoca entre la cantidad y el precio (cantidad que los consumidores están dispuestos a comprar en función del precio del producto, ceteris paribus) y proviene del problema de maximización de utilidad del consumidor [\authorfont{Ver Tema 3.A.8} y \authorfont{Tema 3.A.16}], que presenta una función de utilidad cuasilineal (con efectos cruzados nulos) sujeto a una restricción presupuestaria.\footnote{%
        En 1592, LUIS DE MOLINA economista español de la Escuela de Salamanca, hace uso de la cláusula ceteris paribus: “Cuanto menor es la cantidad de dinero en un sitio, más aumenta su valor y, por tanto, ceteris paribus, con la misma cantidad de dinero se pueden comprar más cosas.” 
    
    Posteriormente, su uso fue popularizado por ALFRED MARSHALL, en su enfoque de equilibrio parcial. LUIS DE MOLINA lo usó para hablar de la teoría cuantitativa del dinero. Se puede argumentar que otro economista español de la misma escuela, MARTÍN DE AZPILCUETA (1556), propuso esta teoría 13 años antes de que lo hiciera el economista francés JEAN BODIN. Concretamente, escribió lo siguiente:
    
    “En las tierras donde hay gran falta de dinero, todas las otras cosas vendibles, y aun las manos y trabajo de los hombres se dan por menos dinero que donde hay abundancia de él; como por la experiencia se ve que en Francia, donde hay menos dinero que en España, valen mucho menos el pan, el vino, los paños, las manos y los trabajos; y aun en España, cuando había menos dinero, por mucho menos se daban las cosas vendibles, las manos y los trabajos de los hombres, que después de que las Indias fueran descubiertas, la cubrieron de oro y plata. La causa es que el dinero vale más donde y cuando falta que donde y cuando es abundante.”
    }\textsuperscript{,}\footnote{%
        Si se dieran efectos cruzados, el poder de mercado se perdería (ya que los agentes percibirían otros bienes como sustitutivos). Este elemento es imprescindible tanto para la agregación de demandas individuales (para construir la demanda agregada de mercado), donde se requiere la existencia de efecto renta nulo, como para el poder del monopolista, donde se requiere, además, efecto-precio-cruzado nulo.
    } 
    \item \textbf{La función de costes totales} ($C(Q)$). Refleja un proceso de optimización previo respecto al empleo de factores de forma que se garantiza eficiencia técnica y económica [\authorfont{Ver Tema 3.A.12}]. Es decir, al utilizar la función de costes ya se supone la eficiencia económica: que la empresa iguala la ${RMST}_{z_j}^{z_i}$ al cociente de los precios de los factores productivos;
    \item \textbf{Restricción de Participación}. Se trata, en puridad, de una restricción de participación. Al permitir la producción nula de la empresa garantizamos que cualquier nivel de producción que la empresa elija proporciona unos beneficios mayores a los de no producir: $\Pi_M(0) = 0$.
\end{la}

\paragraph{Condiciones necesarias: CPO}
La condición de primer orden que nos determina la solución nos indica que el equilibrio viene dado por la igualdad entre ingreso marginal y coste marginal, $IMg = CMg$.\footnote{%
    La condición de igualdad se obtiene para todas las estructuras de mercado y es la solución del problema de maximización del beneficio en dos etapas. COURNOT (1838) fue pionero en llegar a esta conclusión.
}  Esto es, la última unidad producida debe aportar un ingreso igual a su coste.
 
\textbf{NOTA.-} No debemos caer en el error de que en competencia perfecta no se cumple esta igualdad, ya que también se cumple que $IMg=CMg$: en competencia perfecta el $IMg$ es igual al precio, mientras que en un monopolio el $IMg$ es menor que el precio.\footnote{
Recordemos que en competencia perfecta, al cumplirse que el precio es igual al coste marginal, se definía una función inversa de oferta. Sin embargo, en el caso del monopolio tal relación biunívoca no existe porque existen más variables que determinan la formación de precios que el coste marginal. Esto implica que en monopolio la valoración privada del bien (${IMg}_x(Q)$) es menor que la valoración social de dicho bien ($P_x^d$). De manera comparada:
\begin{la}
    \item Competencia perfecta. El productor debe hacer frente a una demanda perfectamente elástica. Es decir, un aumento infinitesimal del precio le hace perder todas sus ventas), lo que en la fórmula antes vista permite que $IMg=P$. 
    
    Es decir, a nivel individual (elasticidad-demanda-precio del productor), si este aumenta los precios, como la elasticidad a título individual es infinita (ya que el equilibrio se mantendrá en el equilibrio competitivo) el $IMg$ será igual al precio. En cambio en monopolio, el monopolista marca las reglas del equilibrio, por lo que la elasticidad-demanda-precio individual no será perfectamente elástica: será negativa (como es obvio), pero finita.
    \item Monopolio. Podemos asegurar, no sólo que la elasticidad será finita, sino que, además, será mayor a uno. Es decir, el monopolista producirá siempre en el tramo elástico de la demanda. 
    
    Se puede demostrar que el monopolista siempre decidirá producir una cantidad situada en el tramo elástico de la función de demanda. Si se situase en el tramo inelástico, siempre podría aumentar hasta cierto punto el precio y reducir la cantidad producida, de tal manera que la reducción de costes superase la caída de ingresos en el tramo elástico. Aumentar los precios reduce los ingresos totales. Pero reduce también los costes, por lo que la empresa tiene margen para reducir la cantidad hasta que ingrese y coste marginal se igualen. Y es que si la demanda es inelástica, el ingreso marginal sería negativo, por lo que no se podría igualar al coste marginal, que es siempre positivo. Visto de otra manera: en el tramo inelástico, una reducción de la producción aumenta el ingreso y, obviamente, disminuye el coste, por lo que aumentaría el beneficio, con lo cual en este tramo no se puede maximizar beneficios. 
    
    Por la misma razón, el monopolista tampoco producirá en el tramo en el que $$|\varepsilon_{D,P}|=1$$ (i.e. $IMg=0$, donde se maximizan los ingresos), porque en ese tramo una reducción de la producción reduce el coste y el ingreso no varía, por lo que el beneficio aumentaría. La única situación en la que produce en $|\varepsilon_{D,P}|=1$ es cuando el $CMg$ es nulo.
\end{la}
} Sin embargo, la diferencia principal radica en que no existe una curva de oferta definida en un mercado monopolístico. A diferencia de competencia perfecta, donde la curva de $CMg$ es la curva de oferta de la empresa, la asociación única entre precios y cantidades no se da en el monopolio, ya que en éste el precio es una variable endógena para el monopolista (lo veremos más adelante). 

\paragraph{Condiciones suficientes: CSO}
Por otra parte, tenemos las condiciones suficientes para el equilibrio (CSO). La CSO, que debe cumplirse para que el beneficio obtenido sea un máximo y no un mínimo, impone que la pendiente del $IMg$ sea menor que la del $CMg$.\footnote{%
    Sin valor absoluto. La pendiente del ingreso marginal será negativa, pero deberá ser más negativa que la del coste marginal (que podría ser negativa o positiva).
} El monopolista, pues, puede maximizar tanto en el tramo creciente de los costes marginales (como exigía la CSO en competencia perfecta) como en el tramo decreciente (siempre y cuando, en este caso, sea más plana que la curva del IMg).\footnote{%
    La condición de segundo orden se cumple automáticamente cuando se está en zona de CMg crecientes. En caso de CMg decrecientes, sería necesario que la pendiente de la curva de CMg fuera menor que la pendiente de la curva de IMg (sin valor absoluto). Si se supone que la función de beneficios es cuasicóncava se garantiza que la cantidad que maximiza el beneficio es única.
}

\paragraph{Índice de Lerner}
Como decíamos, no existe una única curva de oferta del monopolista que pueda derivarse de su curva de $CMg$. Pero, ¿qué significa esto? ¿por qué es relevante? Sabemos que, en competencia perfecta, un incremento infinitesimal del precio por encima del coste marginal implica la pérdida total de la demanda residual de la empresa. Por el contrario, la curva de costes marginales en caso de monopolio nada influye en la determinación del precio ya que, por ejemplo:
\begin{la}
    \item La misma cantidad puede ofrecerse a distintos precios; o,
    \item Se pueden ofertar diferentes cantidades al mismo precio.
\end{la}
	
Ilustremoslo con un ejemplo. Imaginemos dos mercados en situación monopolística, uno es el mercado de insecticidas agrícolas y otro el de una sala de cine en un pueblo. El poder de mercado que ostentan estas empresas es diferente, pero ¿por qué? La empresa de insecticidas tendrá mayor control sobre el precio ya que es un producto cuasi-imprescindible para los agricultores, mientras que una proyección de cine es un servicio relativamente prescindible por parte de los habitantes del pueblo, por lo que el gerente tendrá menos control sobre el precio. La intuición subyacente es la clave de este mercado: la demanda es la que ostenta un papel relevante en la fijación del precio. Y, en particular, su elasticidad-precio. Por tanto, ¿qué es entonces lo que determina la oferta y el poder de mercado del monopolista? Precisamente la estructura de preferencias que subyace en la curva de demanda del mercado. Veamos de nuevo a las CPO para ilustrarlo:

\eqblock{
\begin{aligned}
    &IMg(Q)\equiv \frac{\partial}{\partial Q}P_x^d(Q)\cdot Q+P_x^d(Q)\;\Rightarrow\;\underbrace{\frac{\partial P_x^d(Q)}{\partial Q}\cdot\frac{Q}{P_x^d(Q)}}_{\equiv\frac{1}{\varepsilon_{D,p_x}}}\cdot P_x^d(Q)+P_x^d(Q)\\[1em]
    &\overset{(\varepsilon_{D,p_x}<0)}{\Rightarrow}-\frac{P_x^d(Q)}{|\varepsilon_{D,p_x}|}+P_x^d(Q)\;\Rightarrow\;\boxed{P_x^d(Q)\cdot \left(1-\overbrace{\frac{1}{|\varepsilon_{D,p_x}|}}^{\text{\scriptsize Índice de LERNER}}\right)<P_x^d(Q)}
\end{aligned}
}{Índice de LERNER. Teoría del Monopolio: Monopolio uniproducto}

Como al maximizar estará fijando simultáneamente precio y cantidad, la noción de curva de oferta deja de ser operativa: no existe una relación unívoca entre precio y cantidad basada únicamente en los costes. En consecuencia, la oferta del monopolista queda plenamente caracterizada por la condición $IMg(Q)=CMg(Q)$. 

Si desarrollamos el ${IMg}_x$, obtenemos la desigualdad anterior, que nos indica que el ${IMg}_x$ es estrictamente menor que el precio siempre que la demanda tenga elasticidad finita.\footnote{
$\frac{\partial}{\partial x}U(x)<0$ implica que $P_x^d(x)$ tenga pendiente negativa y, en consecuencia, elasticidad finita: $\varepsilon_{D,p_x}\in(-\infty, 0]$. Si $P_x^d(x)$ presentara elasticidad infinita (demanda horizontal): 
$$P_x^d(x)=\text{cte.}\Longrightarrow {IMg}_x\equiv \underbrace{\underbrace{\frac{\partial}{\partial Q}P_x^d(Q)}_{=0}\cdot Q}_{=0}+P_x^d(Q)\Longrightarrow {IMg}_x(Q)=P_x^d(Q)=\text{cte.}$$ 
Y se alcanzaría la solución competitiva: $IMg_x(Q)=CMg_x(Q)\Longrightarrow P_x^d(Q)={CMg}_x(Q)$.
} Esto ocurre porque, \textbf{al aumentar la cantidad vendida}, \textbf{el monopolista debe reducir el precio} no solo de la unidad adicional, sino \textbf{de todas las unidades vendidas}; y esta\textbf{pérdida inframarginal} es mayor cuanto menos elástica es la demanda. 

El término $\frac{1}{|\varepsilon_{D,p_x}|}$, conocido como Índice de LERNER, mide exactamente cuánto puede separarse el precio del coste marginal sin perder el incentivo a producir. Es decir, el punto en que el ingreso marginal deje de cubrir el coste marginal.\footnote{%
    Si el precio se separa demasiado del coste marginal, entonces: una pequeña expansión de la cantidad provoca una gran caída del precio, el ingreso marginal cae por debajo del coste marginal y producir más deja de ser rentable.
} Entonces, el incentivo a producir está limitado por la elasticidad de la demanda, no por los costes, y: (i) si la demanda es muy elástica, este término es pequeño y el precio se aproxima al coste marginal, limitando el poder de mercado; pero, (ii) si la demanda es poco elástica, el índice aumenta y el monopolista puede fijar precios significativamente superiores al coste marginal.

Así, la desigualdad nos muestra una de las piedras angulares de la Teoría del Monopolio: la importancia del Índice de LERNER, que se define como el margen relativo (\textit{mark-up}) sobre el ${CMg}_x$ y es igual a la inversa de la elasticidad demanda precio.\footnote{%
    Las autoridades de la competencia utilizan este índice para medir el poder de mercado. El nombre de este índice se debe a ABBA LERNER
}

\subsection{Equilibrio: Monopolio en el Mercado}
Una vez conocemos los fundamentos de la Teoría del Monopolio, cabe analizar las propiedades del equilibrio y, en particular, realizar un análisis comparado respecto al mercado de referencia: competencia perfecta.

\subsubsection{Análisis Positivo}
En equilibrio, la empresa monopolista produce hasta igualar el ingreso marginal a los costes marginales y se fija el precio de acuerdo con la curva de demanda. A diferencia de la situación de competencia perfecta:
\begin{la}
    \item El monopolista podría obtener beneficios extraordinarios en el largo plazo. Esto ocurre porque, a pesar de que los beneficios positivos incentivan la entrada de empresas en el mercado, existen barreras a la entrada que impiden que nuevas empresas entren en el mercado, a diferencia de lo que ocurre con la competencia perfecta;
    \item La maximización de beneficios no la realiza el monopolista necesariamente en el mínimo de la curva de costes medios.\footnote{%
        Ejemplo en el que la maximización de beneficios coincide con la escala óptima [gráfico 1]. La condición de maximización de beneficios, IMg = CMg, se produce cuando el CMe alcanza el mínimo. En él se produce un pleno aprovechamiento de la capacidad y de las economías de escala. Ejemplo en el que la maximización de beneficios se produce con una escala subóptima [gráfico 2]. La condición de maximización de beneficios se produce en el tramo decreciente del CMe. En él se da un exceso de capacidad y un aprovechamiento incompleto de las economías de escala. Se trataría de un mercado pequeño que no permite un aprovechamiento completo de la capacidad. Ejemplo en el que la maximización de beneficios se produce con una escala mayor a la óptima [gráfico 3]. La condición de maximización de beneficios se produce en el tramo creciente del CMe. En él se da una sobreutilización y deseconomías de escala. Se trata de un mercado muy grande que obliga al monopolista a producir más allá de su escala óptima. 
    \imagenfit{MOnopolio_AnalisisPositivo_Eq.png}{Análisis positivo del Equilibrio en mercado monopolístico}{Apuntes Víctor Gutiérrez. Tema 3.A.17}
    }
\end{la}

\subsubsection{Análisis normativo}
Desde un punto de vista normativo, debemos comparar los resultados obtenidos con los del modelo de competencia perfecta [\authorfont{Ver Tema 3.A.16}].

\paragraph{Mercado: Pérdidas Irrecuperables de Eficiencia y apropiación}
Ligada a la distorsión en los precios que se produce con respecto al modelo de competencia, existe una pérdida de eficiencia derivada de un mayor precio y una menor cantidad. Esta pérdida se puede cuantificar calculando el area que se correspondería con las áreas $A + B$ del gráfico siguiente:

\imagenfit{PIE_MonopolioNormativo_Eq.png}{}{Apuntes Víctor Gutiérrez. Tema 3.A.17}

Donde:
\begin{la}
    \item \textbf{Consumidor}. El área $A$ representa una pérdida irrecuperable de eficiencia por el lado del consumo, que en competencia perfecta formaría parte del excedente del consumidor; y,
    \item \textbf{Productor}. El área $B$ representa una pérdida irrecuperable de eficiencia por el lado de la producción, que en competencia perfecta formaría parte del excedente del productor.
\end{la}
	
Además, una notable característica es que el área C representa la transferencia de excedentes del consumidor al productor respecto a la situación de competencia perfecta: el \textbf{beneficio extraordinario.}

\paragraph{Mercado de factores: Ineficiecia asignativa}
El grado de ineficiencia en la asignación de recursos puede comprobarse también en lo que respecta a la demanda de factores productivos. 

Suponiendo que los precios de los inputs vienen dados, monopolista precio-aceptante en el mercado de inputs, con precio del factor productivo $\vec{z}=k$. Tenemos que:

\eqblock{
\begin{aligned}
    &\Pi_M(Q)=P_x^d(Q)\cdot Q-k\vec z\qquad \text{donde,}\;\;Q=F(\vec z)\\[1em]
    &\text{Óptimo:}\qquad \frac{\partial}{\partial \vec z}\Pi_M={IMg}_x(Q)\cdot {PMg}_{\vec z}-k=0\\[1em]
    &\underset{\text{\scriptsize Aunque siempre paga $k$}}{\text{Mercado de factores:}}\quad {IMg}_x(Q)\cdot {PMg}_{\vec z}<P_x^d(Q)\cdot {PMg}_{\vec z}
\end{aligned}
}{Ineficiencia asignativa de factores productivos. Teoría del Monopolio} 

Por tanto el monopolista contrata factores hasta el punto en que su ingreso marginal generado, no su valor de producción, iguala el precio del factor. Como el ingreso marginal es menor que el precio, esto implica que el monopolista: (i) demanda menos cantidad de factores; y, (ii) el producto marginal del factor en el óptimo es mayor que en competencia (infrautilización y \textit{sobreremuneración}).

Es importante distinguir esta situación con el caso en el que el monopolista es además monopsonista (único comprador) en el mercado de factores productivos, ya que en ese caso, el precio de los inputs dependerá de la cantidad demandada (como veremos más adelante). En este caso hablaríamos de explotación monopolística de los factores productivos.

\paragraph{Largo plazo: Ineficiente técnica/productiva}
Por último, no se garantiza que la empresa alcance el punto mínimo de la curva de costes medios a largo plazo (conocido como tamaño óptimo de planta escala eficiente de producción).

Bajo monopolio, la elección del tamaño de planta y el grado de utilización de la capacidad productiva no están disciplinados por la competencia, sino que responden exclusivamente al problema de maximización de beneficios del monopolista, condicionado por la demanda del mercado. Dado que el monopolista no tiene incentivo a minimizar costes medios, sino a maximizar beneficios, no se garantiza que opere en la escala óptima eficiente. En consecuencia, puede producir a una escala inferior a la eficiente, manteniendo capacidad ociosa y sin agotar las economías de escala, o bien operar a una escala superior a la óptima desde el punto de vista tecnológico. La ausencia de presión competitiva elimina el mecanismo que, en competencia perfecta, fuerza a las empresas a adoptar la escala eficiente de producción.

Además, la ausencia de competencia reduce los incentivos del monopolista a reducir las ineficiencias implícitas, por ejemplo, la \textit{Ineficiencia X} de LIEBSTEIN.


\subsubsection{Implicaciones de Política Económica}
El debate, tanto teórico como empírico, respecto de las consecuencias de estructuras de mercado monopolísticas ha sido muy intenso.

Desde un punto de vista teórico, cabe distinguir dos posiciones básicas antagónicas: 
	Por una parte, quienes sostienen el análisis de las pérdidas de eficiencia en los términos que se acaban de exponer y que, en general, han aproximado dichas pérdidas para el caso de funciones de demanda lineales, donde es sencillo demostrar que la pérdida neta (deadweight loss) puede aproximarse por la mitad del beneficio del monopolista; 

    
	Por otra parte, el análisis debido a WILLIAMSON (1968) supone que, en general, el sostenimiento de un monopolio se debe a la existencia de ventajas técnicas en forma de economías de escala para el monopolista. Este caso puede venir representado por el hecho de que el monopolio produzca con la tecnología representada por la función de costes etiquetada con «2» del gráfico  (que es menor que bajo competencia- perfecta, C 1 '), en cuyo caso el aumento de excedente del monopolista puede incluso llegar a ser mayor que la pérdida de excedente de los consumidores, dependiendo por tanto la valoración final de la comparación entre las áreas rayadas del gráfico. Los defensores de esta última explicación consideran que la reducción de costes medios de producción por la existencia de monopolio puede estimarse en un 5\%, con lo que las pérdidas por existencia de poder de monopolio tendrían un valor real muy moderado.
 
	Desde un punto de vista empírico, las estimaciones de la pérdida neta de bienestar derivada del monopolio han sido pequeñas. La metodología se basa principalmente en estimar la curva de demanda y la curva de coste marginal. El estudio pionero se debe a HARBERGER (1954), quien estimó una pérdida total de bienestar no superior al 0,1\% de la Renta nacional bruta. Ello dio pie a mucha controversia en torno a los datos y la metodología utilizados por HABERGER.  COWUNG y MUELLER (1978) tratan de tener en cuenta los costes ligados a adquirir y mantener el poder de mercado (fuente de otro efecto perjudicial del monopolio conocido como \textit{rent-seeking}), obteniendo unas estimaciones algo mayores.\footnote{%
	    En TIROLE (1988) se detallan algunas críticas vertidas al enfoque empleado por HABERGER y se analizan más detenidamente estas otras distorsiones, ya sea ligadas a costes (ineficiencia-X) como el comportamiento de búsqueda de rentas (rent-seeking).
	}\textsuperscript{,}\footnote{%
        Cabe resaltar que las pérdidas de eficiencia comentadas son de carácter estático. En la conclusión del tema se realizan algunas consideraciones acerca del monopolio desde una perspectiva dinámica
    }



\clearpage
\section{Discriminación de precios}
\puntosclave{
    El estudio del monopolio multiproducto (esto es, un mismo monopolista, produce conjuntamente distintos bienes en régimen de monopolio) depende crucialmente de las relaciones que se establezcan entre las demandas y costes de producción de los bienes que se producen conjuntamente.
	
    La discriminación de precios se puede definir como la venta de diferentes unidades de un bien a precios distintos, ya sea el mismo consumidor o a distintos, sin que las diferencias en el precio se hayan justificadas por diferencias en los costes de producción. Con la discriminación de precios, el monopolista persigue extraer una mayor parte del excedente del consumidor que bajo un precio uniforme (traslandándolos a su Beneficio).
	
    La clasificación clásica de los tipos de discriminación (debida a PIGOU) se basa en la cantidad de información de la que dispone el monopolista sobre las preferencias de los consumidores:
	
    En la discriminación de primer grado o perfecta, el monopolista conoce la valoración que tiene cada consumidor de cada unidad del producto, pudiendo fijar los precios de modo acorde, por ello conduce al resultado competitivo eficiente.
	
    En la discriminación de segundo grado, el monopolista solo conoce la distribución de las valoraciones de los consumidores, por lo que ofrece el mismo precio a todos los consumidores, pero los precios son distintos para diferentes unidades de producto (asimilándose a una situación de selección adversa).
	
    En la discriminación de tercer grado, el monopolista solo dispone de una señal de las valoraciones de los consumidores, fijando precios distintos en función del grupo al que adscribe al consumidor. En este caso, el problema del monopolista se asemeja al problema de un monopolio multi producto con demandas independientes y costes dependientes.
    
    Posteriormente, se desarrollo la discriminación de cuarto grado o intertemporal, que consiste en establecer diferentes precios en el tiempo a bienes duraderos. El éxito de la misma depende crucialmente de que los consumidores sean impacientes (COASE).

    El rasgo definitorio de un monopolio natural es la existencia de una tecnología de producción que propicie la subaditividad de los costes y la existencia de rendimientos crecientes a escala es una condición suficiente, pero no necesaria del monopolio natural.. El ejemplo paradigmático de monopolio natural son los sectores de red.
    
	Se dice que el monopolio natural es sostenible cuando no se produce la entrada de un nuevo competidor en el mercado. Bajo determinados supuestos (en especial, que el monopolista esté obligado a satisfacer la cantidad total demandada al precio que fije y no pueda modificarlo en caso de entrada de competidores), el monopolio natural solo es sostenible si la tecnología presenta rendimientos crecientes a escala. Si se relajan los supuestos anteriores, el monopolio natural puede ser sostenible bajo rendimientos decrecientes.
}
	
Ahora que hemos visto el marco de la Teoría del Monopolio sabemos cuáles son los fundamentos sobre los cuales opera el monopolio. Una de las características esenciales, como hemos visto, es la inexistencia de la curva de oferta de mercado, que permite al monopolista determinar precios y cantidades a discreción con el fin de maximizar sus beneficios. Ahora bien, esto tiene una consecuencia práctica muy importante: la capacidad de segmentación que se le ofrece al monopolista dentro de su propio mercado, lo que se conoce como \textbf{mecanismos de discriminación de precios}.

Si el poder del monopolista sobre su mercado depende crucialmente de la estructura de preferencias subyacente proyectada sobre la pendiente de la función de demanda de mercado. La capacidad de segmentación depende crucialmente de la información que tenga el monopolista sobre la estructura de preferencias de los consumidores del mercado. Por ello, no es de extrañar que la primera categorización de los mecanismos de discriminación realizada por PIGOU (1920), sucesor de MARSHALL en la Universidad de Cambridge, los distinguiese en función de la cantidad de información de la que dispone el monopolista:
\begin{la}
    \item \textbf{Discriminación de primer grado o perfecta}. Si el monopolista conoce la valoración de cada consumidor el óptimo se alcanzará extrayendo de cada uno su precio de reserva (utilidad marginal individual);
    \item \textbf{Discriminación de segundo grado}. El vendedor sólo conoce la distribución de las valoraciones de los consumidores. En este caso, el monopolista ofrece el mismo precio a todos los consumidores, pero fija distintos precios para diferentes unidades de producto;
    \item \textbf{Discriminación de tercer grado}. El vendedor tiene todavía menos información, sólo una señal de las valoraciones de los consumidores. Por lo tanto, en este caso, el monopolista fija un mismo precio para cada unidad de producto vendida, pero distinto en función del tipo de consumidor.
\end{la}
	
Posteriormente, se ha añadido otro tipo de discriminación a la clasificación clásica de PIGOU, conocida como discriminación de cuarto grado o intertemporal, que consiste en establecer diferentes precios en el tiempo (como puede observarse en los productos tecnológicos, que tienen un precio que va decayendo tras su lanzamiento).

A lo largo de este apartado presentaremos los diferentes tipos de discriminación con el uso de herramientas gráficas (obviaremos el desarrollo analítico, que puede verse en los \textcolor{red}{\textbf{ANEXOS}}). Además, se recomienda recurrir a un ejemplo práctico para hacer más ameno este apartado. 

\begin{specialblock}{\textbf{Ejemplos.-} Mercados Monopolísticos típicos}
    
\end{specialblock}

\subsection{Discriminación de primer grado: Información perfecta}
La discriminación de primer grado es intuitivamente sencilla de comprender: si a cada consumidor le proporciona una utilidad diferente el consumo de un bien que, en última instancia determina su disposición a pagar igualándose con su utilidad marginal, y esta información es conocida por el monopolista; entonces el monopolista fijará un precio $p_{x}=\frac{\partial}{\partial x}U_i(x)$ para cada consumidor igual a su disposición a pagar por la última cantidad que está dispuesto a consumir. Esto da lugar a una segmentación perfecta del mercado que, en términos agregados, conduce a un resultado (en términos positivos) igual al del equilibrio competitivo: la valoración marginal social de la última cantidad intercambiada en el mercado es igual al coste marginal social de la última cantidad producida en el mercado, $P^d_x(x)={CMg}_x$. 

$$\begin{aligned}
    &\text{Equilibrio:}\\[1em]
    &\vec X=(x_1,x_2,\dots,x_I, Q)\qquad \text{cumple}, \quad 
    \left\{\begin{aligned}
        &\underset{\substack{\\[0.5em]\\\text{Para vector de precios:}\;\vec p_x=(p_1,p_2,\dots, p_I)}}{P_x^d(x)=P_x^s(x)\;\sim\;\frac{\partial}{\partial x}{B}_x^S=\frac{\partial}{\partial x}C_x^S}\\[0.5em]
        &\sum_{i=1}^Ix_i=Q
    \end{aligned}\right.
\end{aligned}$$

Gráficamente, el equilibrio se puede representar como:

\imagenfit{PrimerGrado_Eq.png}{Equilibrio: Monopolio discriminador de primer grado}{Apuntes Víctor Gutiérrez. Tema 3.A.17}

En definitiva, la discriminación de precios de primer grado requiere conocer de forma individualizada la disposición marginal a pagar de cada consumidor para cada unidad consumida. Esto equivale a conocer la función de demanda individual completa, es decir, la valoración subjetiva que cada agente asigna a cada unidad adicional del bien. Dado ese conocimiento perfecto de las preferencias, el monopolista puede fijar un precio distinto para cada unidad y para cada consumidor, extrayendo todo el excedente del consumidor y reproduciendo la asignación eficiente desde el punto de vista del output, aunque con una distribución del excedente completamente asimétrica.
 
\subsubsection{Implicaciones de Política Económica}
En términos de Política Económica, este resultado representa un óptimo de PARETO (no mejorable en términos de eficiencia) ya que el vector de equilibrio de mercado $\vec X=(x_1,x_2,\dots,x_I, Q)$ es equivalente al del EGC, que de acuerdo con el 1TFEB es un Óptimo de Pareto.

Por lo tanto, la intervención estatal no estará justificada por motivos de eficiencia y sólo quedará justificada por motivos de equidad. 

\textbf{NOTA.-} Cobrar una tarifa ex-ante de entrada al mercado y luego el precio ($p_x^*$) igual al coste marginal, el resultado sería el mismo. Es decir, los esquemas tarifarios de discriminación perfecta serían equivalentes a la discriminación de primer grado (precios no-lineales), en el sentido de que el monopolista se apropia de la totalidad de los excedentes de cada consumidor. 

\subsubsection{Evidencia empírica}
Es difícil encontrar ejemplos en la práctica, pero cabe citar los zocos o el médico de un pueblo (que cobra a cada paciente lo que esté dispuesto a pagar). Además, algunos de los siguientes artículos han sido publicados al respecto:

Ejemplo práctico en la industria hotelera y de aerolíneas, donde sistemas de \textbf{revenue management} se aproximan a la discriminación de primer grado mediante precios personalizados en tiempo real. Por otro lado, precios personalizados basados en el historial de compras del consumidor, representando una forma de discriminación de primer grado facilitada por datos digitales.\footnote{%
    PHILLIPS (2005), “Pricing and Revenue Optimization”
ACQUISITI y VARIAN (2005), “Conditioning Prices on Purchase History”
}

La discriminación de primer grado resulta evidentemente difícil ponerla en práctica ya que necesitamos información muy detallada sobre los agentes de mercado. Aunque esto es cada vez más sencillo gracias a las herramientas informáticas y el procesamiento de grandes volúmenes de datos existen formas alternativas más sencillas.

\subsection{Discriminación de segundo grado: Auto-selección}
Lo expuesto resulta excesivamente exigente en términos de información pero es lógico suponer que, con una trayectoria suficientemente dilatada en el tiempo, el monopolista es consciente de la existencia de una heterogeneidad de disposiciones marginales a pagar por cada unidad consumida y puede recurrir a técnicas menos exigentes en información para perseguir sus objetivos. 

La más extendida es la discriminación de segundo grado: el monopolista no observa los tipos individuales ni las valoraciones subjetivas de cada consumidor, pero sí conoce la existencia de diferentes grupos de consumidores que presentan distintas disposiciones marginales a pagar para cada unidad adicional consumida. En este contexto, el monopolista puede inducir la auto–selección mediante menús de contratos óptimamente diseñados. Estos menús o esquemas precio–cantidad inducen la revelación del tipo a través del \textbf{patrón de consumo} seleccionado (se conoce también como \textbf{intensidad del consumo}). Por tanto, la segmentación del mercado se consigue mediante la fijación de diferentes esquemas, toda vez que evita que los consumidores con una determinada estructura de preferencias imiten el comportamiento de consumidores con otra estructuras de preferencias diferentes a la que se ha asignado un esquema o menú diferente. 

Consideremos un mercado monopolístico en el que el monopolistia ha identificado dos grupos de agentes diferente: un grupo que presenta alta valoración marginal (demanda más inelástica) y un grupo que presenta baja valoración marginal (demanda más elástica). Gráficamente:



Podemos distinguir dos grandes familias de mecanismos según el momento y la forma en que se extrae el excedente del consumidor. La diferencia fundamental reside en si la revelación del tipo se produce antes del consumo o a través del patrón de consumo mismo, y en qué instrumento contractual se utiliza para inducirla.

\subsubsection{Esquemas de precios}
El primer grupo corresponde a los mecanismos de discriminación basados en auto–selección vía cantidad o calidad, que pueden denominarse esquemas de discriminación por consumo o revelación ex post. El monopolista ofrece un menú de precios–cantidad, bundles o versiones de calidad diferenciada, y la revelación del tipo se produce a través de la elección efectiva del consumidor, una vez observado el esquema de precios. Estos mecanismos explotan diferencias en la valoración marginal y en la intensidad de consumo, no en la disposición total a pagar. Por tanto, el excedente se extrae de forma inframarginal y progresiva, a medida que el consumidor consume más unidades o elige versiones de mayor calidad.

\paragraph{Descuentos por cantidad: \textit{Rappel} sobre ventas}
El tipo más extendido de este grupo es el descuento por cantidades: \textit{rappel sobre ventas}. ¿Quién no se ha enfrentado alguna vez a un descuento por volumen?

A priori, puede resultar una estrategia empresarial que busca perseguir otros fines más allá de la maximización del beneficio: confianza, fondo de comercio, reducción de existencias, etc. pero no pierde en esencia su objetivo discriminador. De hecho, una interpretación ligera y superficial nos conduciría a pensar que el grupo de consumidores a los que va dirigido dicho esquema salen beneficiados por él. 

No obstante, aunque es \textit{cierto}, el razonamiento es más bien inverso: este grupo de consumidores tiene una elevada valoración marginal por lo que, al ofrecer el descuento inframarginal sobre cantidad consumida, estos aumentarán su consumo y el monopolista será capaz de extraer excedente adicional que, aunque en términos individuales sea inframarginal, en términos agregados respresenta un aumento considerable de la apropiación del excedente. Por otro lado, los consumidores con menor valoración marginal y demanda más elástica limitarán su consumo a los tramos iniciales del menú, que ya consumían antes. De este modo, el esquema de descuentos por cantidad permite al monopolista segmentar la demanda sin observar directamente los tipos, explotando las diferencias en valoración marginal a través de la auto–selección.

La clave está en qué excedente es explotable bajo información imperfecta y qué restricción impone la auto–selección: si el monopolista mantuviera un precio único alto, efectivamente extraería mucho excedente unitario de los consumidores con alta valoración, pero perdería beneficios inframarginales por cada unidad no producida: esos consumidores comprarían menos cantidad de la que sería óptima venderles en caso de segmentación real de mercado. El objetivo, por tanto, no es atraer a los tipos bajos (que aportan poco), sino expandir la cantidad vendida a los tipos altos sin bajar el precio de las primeras unidades vendidas. Gráficamente:



El excedente adicional que el monopolista capta proviene de esas ventas inframarginales adicionales a los tipos altos, no de los tipos bajos.

\paragraph{Empaquetamiento de productos (bundling)}
Otra forma de discimación por auto-selección muy comúnmente utilizada se da cuando el monopolista ofrece varios productos. Bajo la misma óptica que en el problema anterior, un monopolista buscará diseñar descuentos óptimos para un paquete conjunto de bienes, de modo que se ofrecen varios productos a un precio agregado inferior que si se compran cada uno por separado.\footnote{%
    En esencia, supone la extensión del problema básico de discriminación de segundo grado a un contexto multidimensional, por vender el monopolista varios productos.
} Esta práctica se conoce por el nombre genérico de empaquetamiento de productos o bundling y se halla muy extendida en multitud de campos, entre ellos los llamados \textit{information goods} (piensese en Microsoft Office).

\paragraph{Calidad diferenciada}
La discriminación mediante la elección de la calidad del producto (MUSSA y ROSEN, 1978) se pliega a un análisis similar a la de cualquier mecanismo de revelación por auto-selección, en tanto en cuanto se puede asemejar la calidad a la valoración del bien (siempre y cuando no haya diferencias significativas en cuanto al coste de producción). 

Un ejemplo evidente sería la edición de libros de tapa blanda (o rústica) y los de tapa dura.

\begin{specialblock}{\textbf{Ejemplo.-} Paradoja BARBIE}
    Otra ejemplo especialmente significativo es el de la muñeca Barbie, en todos los años que lleva en el mercado, se va disfrazando de diferentes profesiones (p.ej. cocinera, médica, deportista, jardinera, etc.) y el precio de ésta varía enormemente en función de la profesión.  Es decir, la Barbie médica puede costar en torno a 25 € mientras que otras Barbies, como la Barbie jardinera, sólo cuestan en torno a 12 €. Pero, ¿cómo se explica esta diferencia en precios? Fijándonos en los costes, no podemos concluir que el coste de producir la Barbie doctora sea diferente al de producir la Barbie jardinera.  

    La diferencia en precios se debe a la discriminación de precios de tercer grado que está realizando la empresa. Ésta sabe que los objetivos que tienen los padres con respecto a sus hijos son muy diferentes en función de su riqueza y, de este modo, las familias con mayor nivel de renta tendrán mayor disponibilidad a pagar puesto que desean que sus hijos obtengan trabajos mejor remunerados (médicos, abogados, etc.). Por el contrario, las familias con menores recursos, que tendrán menor disponibilidad a pagar, se conformaran con regalar una muñeca cualquiera a sus hijos, independientemente de la profesión    
\end{specialblock}


\subsubsection{Esquemas mixtos: Tarifas}
El segundo grupo corresponde a los mecanismos basados en tarifas no lineales o contratos de acceso, también conocidos como esquemas de dos partes o mecanismos de extracción ex ante. En estos esquemas, el monopolista diseña contratos que separan una cuota fija (o pago de acceso) y un precio por unidad, o directamente una tarifa completamente no lineal. El consumidor revela su tipo antes de consumir, aceptando o no el contrato. Por tanto, la discriminación se realiza en el margen extensivo (entrar o no entrar) y en el pago fijo, mientras que el consumo posterior queda internalizado. Desde el punto de vista informacional, estos mecanismos explotan la heterogeneidad en la disposición total a pagar, más que en el patrón marginal de consumo.

En síntesis, a través de este mecanismo el monopolista tratará de segmentar el mercado mediante menús óptimos de dos componentes esenciales:
\begin{la}
    \item \textbf{Tarifa}. Una tarifa o cuota de enntrada, $A$, que otorga derecho a entrar en el mercado; y,
    \item \textbf{Precio}. Un precio estipulado, bien sea fijo o variable, internalizado. Los consumidores efectivos del mercado son solo aquellos que han \textit{entrado}.
\end{la}

El problema al que se enfrenta el monopolista es determinar la tarifa de entrada ($A$) y la de uso ($p_x$)  sujetas a dos restricciones fundamentales:
\begin{la}
    \item Restricción de participación (IR): cada tipo debe preferir participar antes que no comprar; y,
    \item Restricción de incentivos (IC): cada tipo debe preferir el contrato “pensado para él” frente al contrato del otro tipo.
\end{la}

Lo que distingue a los tres esquemas tarifarios es qué instrumentos contractuales están disponibles y, por tanto, qué grado de extracción de excedente es posible.

\paragraph{\textit{Two-part simple tariff}: Tarifa exclusivista}
En una tarifa exclusivista, el monopolista ofrece un único contrato compuesto por:
\begin{lnum}
    \item \textbf{Tarifa}. Una cuota fija de acceso $A$; y,
    \item \textbf{Precio}. Un precio por unidad $p_x$.
\end{lnum}

El contrato es exclusivista en el sentido de que no hay alternativas dentro del menú: todos los consumidores que compran lo hacen bajo las mismas condiciones.

Dado que el monopolista no puede discriminar entre tipos, debe fijar la tarifa teniendo en cuenta al tipo con menor disposición a pagar, ya que si este no participa, el mercado se pierde parcialmente. En equilibrio, la cuota fija queda limitada por la utilidad de reserva del grupo de baja valoración, lo que implica que el tipo alto obtiene \textbf{rentas informacional positivas}. Por tanto, aunque este esquema permite cierta extracción de excedente, es ineficiente desde el punto de vista discriminatorio, porque el monopolista renuncia a capturar parte del excedente del tipo alto para no excluir al tipo bajo.

\paragraph{\textit{Screening}: Tarifas no exclusivistas}
En una tarifa no exclusivista, el monopolista ofrece un menú de contratos diseñados para inducir \textit{auto-selección}. Por ejemplo, supongamos que el monopolista ofrece dos tipos de contratos:
\begin{la}
    \item \textbf{Básico}. Un contrato básico, dirigido al tipo de baja valoración,mediante el cual satisface exactamente su restricción de participación (utilidad de reserva); y,
    \item \textbf{Premium}. Un contrato premium, dirigido al tipo de alta valoración, mediante el cual satisface su restricción de incentivos (no quiere imitar al tipo bajo), pero no necesariamente su restricción de participación con igualdad.
\end{la}

Por tanto, ambos contratos están disponibles para todos, y el monopolista no impide que un tipo elija el contrato del otro. La separación se logra endógenamente, a través del diseño de cantidades y pagos.

Como consecuencia, el tipo alto obtiene una renta informacional positiva, mientras que el tipo bajo no. Para evitar la imitación, el monopolista distorsiona a la baja la cantidad ofrecida al tipo bajo, alejándola del nivel eficiente. Esta distorsión es el precio que paga el monopolista como renta informativa al agente de alta valoración por las asimetrías informativas.

\paragraph{Perfectamente no-lineales}
El desarrollo más reciente de discriminación de segundo grado es la tarifa totalmente no lineal (TIROLE, 1988). Imaginémos que el monopolista puede fijar una función de pagos completa T(q), no solo un precio por unidad y una cuota fija, y que el consumidor elige libremente la cantidad, pero el pago total depende de forma arbitraria de dicha cantidad.

Este instrumento es extremadamente potente, porque permite al monopolista hacer coincidir elección de cantidad y tipo de valoración de manera casi perfecta pero sin conocer verdaderamente la estructura de preferencias individualizadas. Es decir, se alcanza un resultado muy parecido al de la discriminación perfecta mediante la auto-selección, que se logra a través de la forma de la tarifa aunque persiste la propiedad de no-distorsion at the top (los consumidores de alta valoración mentendrán algo como renta informativa).

\subsubsection{Implicaciones de Política Económica}
Existe multitud de evidencia empírica al respecto. En relación con los primeros tiposde discriminación (tarifas), podemos citar:
	La estimación del incremento de beneficios de la empresa debido a esquemas de precios no-lineales.  
	Multitud de mercados caracterizados porser industrias de red. Entre otros: la telefonía, la distribución de agua, la distribución de electricidad, etc. 
Por otro, lado, respecto a los otros dos tipos de discriminación de segunda grado, podemos incluir:
	La edición de libros rústica bajo encuadernación en rústica (popularmente conocida como tapa blanda) frente a la encuadernación cartoné (tapa dura).
	Bundling, por ejemplo paquetes de software informático - Office-, suscripción a revistas online, etc., y más recientemente, en los servicios de comunicaciones electrónicas (los llamados 'paquetes convergentes').


En definitiva, el monopolista conoce la estructura de la demanda agregada la existencia de heterogeneidad en las valoraciones marginales. Por ello, busca endógenamente la segmentación del mercado por autoselección: ofrece un menú de contratos, tarifas o combinaciones precio–cantidad diseñadas de forma que cada consumidor revele indirectamente su tipo a través de la opción que elige. La discriminación se realiza, por tanto, sobre la cantidad consumida, y el instrumento fundamental es el precio no lineal. 

\subsection{Discriminación de tercer grado: Características observables}
Ahora bien, imaginémos que el monopolista puede observar alguna característica verificable del consumidor y esto le permite segmentar el mercado antes de que el consumidor enfrente ningún problema de decisión (por ejemplo, edad, localización, pertenencia a un colectivo, momento del consumo). Aunque el monopolista siga sin conocer la valoración exacta de cada individuo, puede identificar a qué grupo pertenece y aplicar precios lineales distintos en cada segmento. Por tanto, lo que estaría haciendo es segmentar objetivamente el mercado en función de la/s característica/s observada/s ex ante. 

Esto es lo que se conoce como discriminación detercer grado. En este contexto, la segmentación es exógena y observable para el monopolista y cada submercado tiene su propia curva de demanda agregada. Por tanto, el monopolista resuelve un problema de maximización independiente en cada uno, igualando ingreso marginal y coste marginal.\footnote{%
    Los costes dependientes complicarían el problema. Esta situación es característica de problemas dinámicos, q1 y q2 hacen referencia a cantidades en distintos períodos (no distintos productos). Este es el caso del llamado efecto aprendizaje en el que la producción pretérita permite abaratar los costes en el presente y que podemos encontrar en la industria aeronáutica. Esto lleva a aumentar la producción en el primer periodo y permite reducir el precio en el primer período con respecto al precio de monopolio uniperíodo. Sin embargo, lleva a transferir poder de monopolio al futuro.
} En términos económicos, mientras la discriminación de segundo grado discrimina por intensidad de consumo, esta discrimina por características observables del consumidor (pertenencia a un grupo).

\textbf{AQUÍ ES MUUUUUUY FÁCIL SEGUIR CON EL EJEMPLO DE AEROLINEAS}



Habitualmente, las compañías de transporte aplican tarifas especiales para jóvenes o tienen tarifas business enfocadas para hombres de negocios con mayor disponibilidad a pagar.

\subsubsection{Implicaciones de Política Económica}
A pesar de que comúnmente la discriminación de precios sugiere rechazo (debido al trato desigual), no tiene por qué ser necesariamente negativa en términos de bienestar total. Así, la discriminación de tercer grado puede dar lugar a mejoras en el sentido de PARETO cuando con precios uniformes existe un grupo con baja valoración que no compra porque el precio está fijado como el de monopolio para el grupo de alta valoración, y la discriminación permite la compra del producto al grupo con baja valoración antes excluido. En general, podría aumentar el excedente total, siempre y cuando aumente la cantidad producida.

La discriminación de precios no puede aumentar el excedente total a menos que conduzca a un incremento en la producción total, por lo anto, cabría hacernos la siguiente pregunta: En un monopolio que aplica una discriminación de tercer grado, ¿la eficiencia social mejora o empeora respecto al caso de un monopolio puro? No existe una respuesta inequívoca, puede mejorar o empeorar. Lo que está claro es que el beneficio del monopolista no empeora (aumenta o, por lo menos, se mantiene constante).

\subsection{Discriminación de cuarto grado: Intertemporal}
El último tipo de discriminación que veremos es muy peculiar. Consiste en fijar distintos precios en el tiempo a un bien duradero (i.e., bien que, una vez comprado, se consume durante un período prolongado de tiempo). 

En presencia de un bien duradero, resulta improbable que un consumidor que compra un bien duradero hoy lo compre mañana, por lo que los bienes ofrecidos por el monopolista en dos periodos distintos son sustitutivos en lugar de complementarios. Como consecuencia, el monopolista de un bien duradero crea su propia competencia: al vender hoy, reduce su demanda mañana. 

Pero, ¿esta particularidad es positiva o negativa para el monopolista? Una primera aproximación nos haría pensar que esto juega en contra del monopolista. Nada más lejos de la verdad: la posibilidad de vender un bien a lo largo de una franja dilatada en el tiempo otorga al monopolista la posibilidad de extraer paulativamente una cantidad mayor de excedente del consumidor. ¿Por qué? Precisamente porque el monopolista podrá ejercer el poder de mercado que ostenta de forma segmentada por el tiempo y una cualidad fundamental de los agentes económicos: la impaciencia.\footnote{%
    Siguiendo la literatura económica, esta estrategia de reducir precios con el paso del tiempo ha sido denominada Pacman o comecocos ya que, al reducir el precio del bien duradero con el tiempo desde aquel dirigido para el consumidor que tiene mayor disponibilidad a pagar (para así capturar todo el excedente del consumidor), el monopolista va comiéndose los consumidores a lo largo de la curva de demanda. El éxito de esta estrategia depende, crucialmente, de que los consumidores sean impacientes (i.e. que no puedan retrasar la decisión de compra).
} 

Veámoslo gráficamente:


En este contexto, COASE (1972) enunció la \textbf{conjetura de COASE}, que establece que un monopolista que vende un bien duradero pierde su poder de monopolio si:\footnote{%
    El problema del monopolista se asemeja a un juego dinámico con horizonte finito (por lo que resulta de aplicación el equilibrio de NASH perfecto en subjuegos, y el procedimiento de inducción hacia atrás), cuyo desenlace dependerá de si el anuncio de no rebajar los precios en el futuro resulta creíble.

CHURCH y WARE (2000) recogen una lista de las posibles estrategias que puede seguir un monopolista para escapar de la conjetura de COASE (página 135 y ss.).
}
\begin{la}
    \item \textbf{Frecuencia de \textit{re pricing}}. El aumento de la frecuencia incrementa el incentivo del consumidor a esperar, lo que analíticamente se puede interpretar como que la elasticidad de la demanda tiende progresivamente a infinito, lo cual obliga al monopolio a reducir su precio. 
    \item \textbf{Paciencia}. En cambio, se puede mostrar que, cuando un monopolista logra comprometerse de forma creíble a no modificar sus precios en el tiempo, sus beneficios son mayores a los que obtendría si decidiese luego modificarlos. Por lo que la única alteración posible vendría ocasionada por un cambio en las preferencias temporales de los consumidores.
\end{la} 

\subsubsection{Implicaciones de Política Económica}
Ello explicaría, por ejemplo, la práctica observada por parte de algunos monopolistas de bienes duraderos consistente en adoptar un modelo de leasing en lugar de un modelo puro de venta; o, por otro lado, las estretegias de venta presentes en los productos tecnológicos.







\clearpage
\section{Extensiones a la Teoría del Monopolio}
\puntosclave{
	El monopsonio es aquella situación en la que una empresa es la única demandante de un input. Bajo ciertos supuestos simplificadores, el número óptimo de unidades que comprará el monopolista se determina siguiendo el mismo principio que el monopolio (unicproducto).
	
    En un monopolio bilateral (encuentro entre un monopolista y un monopsonista), las condiciones de costes y demandas son insuficientes para determinar el precio y la cantidad, por lo que hay que acudir para ello a factores ajenos a la teoría tradicional, como la habilidad negociadora, el azar, etc.
}
Hasta aquí hemos visto las características esenciales de un monopolio y como éste puede actuar para explotar en mayor medida la estructura del mercado en el que opera. 

Ahora bien, quedan algunas cosas muy importantes que debemos introducir. Por ejemplo, nada se ha dicho sobre el origen de este etstructura de mercado. También sería interesante analizar cómo debería actual un monopolio que ostenta poder en diferentes mercados para maximizar sus beneficios conjuntos. O., ¿esiste la posibilidad de tener un monopolio inverso? ¿existe el poder de \textit{demanda} de mercado? ¿cómo se concibe? ¿qué implicaciones tiene? ¿qué resulta de la interacción de poder en ambas caras del mercado?

Todo esto son cuestiones que pasaremos ahora a analizar brevemente. En concreto, veremos extensiones que centran su atención en: (i) la interacción del productor monopolístico entre diferente mercados; (ii) la condición de comprador monopolístico, esto es, el monopsonio; y, (iii) la interacción de poder en ambas caras del mercado: monopolio y monopsonio.

\subsection{Monopolio natural}
En la introducción se ha mencionado como posible causa de la existencia de un único productor la presencia de una tecnología de producción que propicie la subaditividad de los costes.  Dicha situación se conoce como monopolio natural . 
La evolución histórica del concepto de monopolio natural, se distingue por haber experimentado dos etapas:
	Existencia debido a rendimientos crecientes a escala. Inicialmente, para modelizar la ventaja absoluta de costes que caracteriza al monopolio natural, se supuso que la producción del bien presenta rendimientos crecientes.
	Alteración dinámica: interrelación entre costes y demanda. Sin embargo, la existencia de rendimientos crecientes sólo constituye una condición suficiente, pero no necesaria, del monopolio natural porque no existen rendimientos crecientes para todo nivel de producción. Tarde o temprano algún factor limitativo hará entrar la producción en zona de rendimientos decrecientes, por lo que contradiría la existencia como consecuencia de RCrE. 
Por tanto, en el análisis del monopolio natural subyacen dos grandes cuestiones relacionadas entre si: (i) la deseabilidad u existencia (que no siempre requiere la existencia de rendimientos crecientes); y, (ii) su sostenibilidad en el tiempo (esto es, el análisis de las condiciones bajo las cuales no se produce la entrada de un nuevo competidor en el monopolio natural). 
4.3.1. Deseabilidad o existencia
El concepto de monopolio natural es siempre un concepto relativo a los costes y a la demanda (hay que tener en cuenta el progreso tecnológico y la evolución de la demanda) , por lo que es preciso una explicación del monopolio natural que no se apoye exclusivamente en los rendimientos de escala.
4.3.1.1. Costes
Suponiendo que la producción del bien puede ser llevada a cabo por tantas empresas como se desee, incurriendo cada una de ellas en un coste de producción $C_i (x)$, los costes presentan subaditividad para un volumen de producción determinado si resulta menos costoso producir dicha cantidad en una sola empresa que entre varias:
 
Gráficamente, la región de subaditividad de los costes es muy intuitiva: se trata del rango de output en el cual la producción monopolística presenta costes inferiores a la producción distribuida entre dos o más empresas. Además, si representamos la curva de CMe de una y dos empresas (en forma de U) veremos como podemos combinar el supuesto de RCreE y el rango de subaditividad de los costes.
 
Sin embargo, de este mismo gráfico vemos que se desprende una condición necesaria para la existencia de un monopolio natural: que la función de demanda agregada también coincida con el rango de subaditividad de costes. Por tanto, pasemos a estudiar esto en más detalle.
4.3.1.2. Demanda
Analizando las diferentes curvas de demanda podemos distinguir entre diferentes situaciones:
	Con la función de demanda $D_1$, existen condiciones de monopolio natural en el sentido de que la producción total se obtiene a costes menores en una sola empresa;
	Si la demanda crece hasta adoptar una posición como $D_2$, situándonos en el tramo de rendimientos decrecientes, sigue siendo técnicamente más eficiente el monopolio, ya que $D_1$ va por debajo de la curva de costes medios agregada de dos empresas (nos encontramos todavía en la zona de subaditividad de los costes).
	Una demanda como $D_3$ representaría el caso de un duopolio natural y sostenible.  
Por este motivo, hoy en día se entiende que la condición de monopolio natural (o socialmente deseable) es la subaditividad de los costes. 

4.3.2. Sostenibilidad
Se dice que el monopolio natural es sostenible si no se produce la entrada de un nuevo competidor en el mercado ya que, que la industria tenga características de monopolio natural no significa que el monopolio sea por sí solo sostenible.
Como ocurría con la deseabilidad o existencia, que el monopolio natural sea sostenible o no va a depender de dónde se sitúe la demanda. En efecto, como señala JULIO SEGURA 
[PONER LOS DOS EJEMPLOS: DEMANDA POR ENCIMA Y DEMANDA POR DEBAJO]
Si el caso comentado es de aplicación general, ante la ausencia de barreras de entrada y salida, la sostenibilidad de la situación de monopolio (de producción simple) exige que la tecnología presente rendimientos crecientes, que es la conclusión sostenida por los defensores de la teoría de los mercados atacables o impugnables (contestable markets).  Sin embargo, la relajación de los supuestos mantenidos puede permitir la sostenibilidad del monopolio incluso bajo rendimientos decrecientes: 
	En primer lugar, las barreras a la entrada no sólo provienen de intervenciones legales o administrativas, sino de comportamientos estratégicos del monopolista instalado; 
	En segundo lugar, considerar que el Incumbente acomodará al Entrante suministrando únicamente a la demanda residual es poco realista porque implica que el monopolista no reaccionará reduciendo sus precios cuando entre un nuevo competidor vendiendo a precios inferiores al suyo, aunque ello le acarree pérdidas muy superiores. En otras palabras, implica suponer que la velocidad de reacción del instalado es menor que la del entrante y que la de los consumidores; 
	Por último, la existencia de costes irrecuperables (sunk costs) actuaría como barrera de salida para los competidores potenciales, que es probable que se vean expulsados del mercado ante una reacción menos pasiva de la empresa ya instalada. En cuyo caso, el monopolista no se vería obligado a vender a un precio igual al coste medio.
Por lo tanto, el concepto de monopolio natural es siempre un concepto relativo a los costes y a la demanda (hay que tener en cuenta factores como el progreso tecnológico y la evolución de la demanda). 
4.3.3. Extensiones: Monopolio natural multiproducto
El concepto clave no sería la subaditividad de costes, sino el de economías de alcance. Para el caso de dos productos, la condición de existencia o deseabilidad del monopolio natural multiproducto viene dada por:

Es decir, que el coste de producción en el seno de una empresa de dos productos sea menor que la suma del coste de dos empresas independientes que produzcan cada una uno de los productos.
4.3.3.1. Sostenibilidad
Un monopolio multiproducto, será sostenible si fija una estructura de precios que no incluye subsidios cruzados. 
Si un monopolio, para poder obtener beneficios, tiene que fijar una estructura de precios que incluya subsidios cruzados, no será sostenible porque un competidor podría producir sólo los bienes que permiten mantener la estructura de subsidios cruzados y venderlos a precios inferiores que el monopolista. Por tanto, para que un monopolio multiproducto sea sostenible tiene que poder fijar precios que no impliquen subsidios cruzados. 


\subsection{Monopolio multiproducto: producción conjunta}
¿Qué pasaría en el caso de que el monopolio no produzca un solo bien, sino varios (monopolio multiproducto)? ¿Puede optimizar el monopolio sus beneficios conjuntos? ¿Debe hacerlo?

Todas las respuestas depende crucialmente de dos elementos: ¿son las demandas en cada mercado independientes o, por el contrario, son dependientes entre sí? Y, ¿con los procesos productivos, y en especial los costes, separables o no? Así distinguiremos los siguientes casos:
\begin{la}
    \item Demandas independientes y costes separables;
    \item Demandas dependientes y costes separables; y,
    \item Demandas independientes y costes no separables.
\end{la}
 
\subsubsection{Demandas independientes y costes separales}
En el caso de monopolio multiproducto con demandas independientes y costes separables, cada bien se tratará como un mercado separado. Por lo tanto, la resolución será idéntica que en el caso de monopolio uniproducto para cada uno de los mercados en los que opere.

\subsubsection{Demandas dependientes y costes separables}
En el escenario de un monopolio multiproducto con demandas dependientes y costes separables, el problema de optimización es análogo al del monopolio uniproducto, con la salvedad de que se incluyen todos los bienes producidos en la función objetivo. El problema se presentaría de la siguiente forma:

\eqblock{
\begin{aligned}
    &\text{Costes separables}\;\Longrightarrow\;C(\vec q)=\underset{\text{\scriptsize No economías de alcance o interacciones productivas}}{\sum_{i=1}^IC_i(q_i)}\\[0.5em]
    &\text{Demandas dependientes}\;\Longrightarrow\;\underset{\equiv\text{\scriptsize  Sustitución o complementariedad entre productos}}{\frac{\partial P_{i}^d(\vec q)}{\partial x_j}\neq0,\quad(i\neq j)}\\[2em]
    &\arg\max_{\{\vec q=(q_1,q_2,\dots,q_I)\}}\Pi(\vec q)=\sum_{i=1}^IP_{i}^d(\vec q)\cdot q_i-\sum_{i=1}^IC_i(q_i)\\[1em]
    &\text{CPO:}\qquad \frac{\partial}{\partial q_k}\Pi(\vec q)=\underbrace{P_{k}^d(\vec q)}_{\text{\scriptsize Ingreso directo}}+\underbrace{\sum_{i=1}^I\frac{\partial P_{i}(\vec q)}{\partial q_k}\cdot q_i}_{\text{\scriptsize Efectos cruzados sobre todos lo productos}}-\underbrace{\frac{\partial}{\partial q_k}C_k(q_k)}_{\text{\scriptsize Coste marginal}}=0\\[0.5em]
    &\underset{\text{\scriptsize reordenando}}{\Longrightarrow}\qquad \underbrace{P_{k}^d(\vec q)+\sum_{i=1}^I\frac{\partial P_{i}(\vec q)}{\partial q_k}\cdot q_i}_{\text{\scriptsize Ingreso marginal multiproducto}}=\frac{\partial}{\partial q_k}C_k(q_k)\;\Longrightarrow\;{IMg}_k=\frac{\partial}{\partial q_k}C_k(q_k)\\[2em]
    &\underset{\text{\scriptsize Índice de LERNER generalizado}}{\text{Regla de fijación de precios:}}\qquad \boxed{\mathbf{m}=-\left(\mathbf{E}^\top\right)^{-1}\mathbf{s}}\quad\text{donde,}
    \left\{\begin{aligned}
        &\mathbf{m}\equiv&&\substack{\text{\scriptsize Vector de margenes relativos}\\\text{\scriptsize de precio-coste marginal}}\\[0.5em]
        &\left(\mathbf{E^\top}\right)^{-1}\equiv&&\substack{\text{\scriptsize Matriz de elasticidades-precio de demanda}\\\text{\scriptsize (propias, diagonal, y cruzadas, resto)}\\\text{\scriptsize transpuesta e invertida}}\\[0.5em]
        &\mathbf{s}\equiv&&\substack{\text{\scriptsize Vector de cuotas de}\\\text{\scriptsize ingreso, mide importancia relativa}\\\text{\scriptsize de los ingresos totales}}\\[0.5em]
    \end{aligned}\right.
\end{aligned}
}{Índice de LERNER general: Monopolio multiproducto. Extensiones: Teoría del Monopolio}

Por tanto, el índice de Lerner incluirá un término adicional (además de la elasticidad precio-propia): la elasticidad precio-cruzada. De esta forma, los bienes podrían ser:
\begin{la}
    \item \textbf{Sustitutivos brutos}. Si la elasticidad-precio-cruzado es mayor que cero. En cuyo caso el Índice de Lerner será mayor al caso del monopolio monoproducto (i.e. mayor que la inversa de la elasticidad-precio propio) y el monopolio fijará precios más altos que si cada bien fuera producido por separado por empresas independientes, dado que tiene la ventaja de que cuando aumenta el precio de uno de los bienes conduce a una expansión de la demanda del otro. Esto ocurre porque el monopolista tiene más dominio sobre los consumidores (estos no pueden escapar); \footnote{%
        Esto podría explicar, por ejemplo, porque un supermercado físico, si pasa a vender online, fijará precios más altos que un supermercado que sólo esté presente en internet.
    } y, 
    \item \textbf{Complementarios brutos}. Si la elasticidad-precio-cruzado es menor que cero, el Índice de Lerner será menor al caso del monopolio monoproducto (i.e. mayor que la inversa de la elasticidad-precio propio) y el monopolio fijará precios más bajos que si cada bien fuera producido por separado por empresas independientes. La idea es que fijar precios más altos llevaría a una pérdida de la demanda de todos los bienes. Incluso puede darse que fuera conveniente en este caso vender algunos bienes a un precio inferior a su coste marginal, de modo que su Índice de Lerner llegaría a ser negativo.
\end{la}
	
\subsubsection{Demandas independientes y costes no separables}
En el caso de monopolio multiproducto con demandas independientes y costes no separables, la demanda de cada producto depende exclusivamente de su precio. Sin embargo, los costes exhiben efectos cruzados no-nulos. Esta estructura se conoce comúnmente como \textit{learning-by-doing} o economías de alcance. Es decir, la producción de un bien permite abaratar la producción del otro.\footnote{%
    En un entorno dinámico, podemos considerar que los bienes que se están tratando son diferentes \textit{modelos} en un intervalo temporal dilatado. En este caso estaríamos diciendo que la producción pretérita permite abaratar los costes en el presente y, por ejemplo, este tipo de análisis es común en la industria aeronáutica.
}

Si el efecto aprendizaje es muy intenso, puede darse el caso en el que se decidiese fijar un precio inferior al ${CMg}_x$ en uno de los bienes, lo que sería un ejemplo de subsidios cruzados. Es decir, aunque la producción de un bien genera pérdidas aisladas, se subsidia en el otro mercado son el fin de abaratar los costes conjuntos de fabricación. 

Como comentarios finales se puede señalar que: por un lado, pueda parecer que no existe poder de monopolio aunque lo que sucede es que se transfiere hacia el futuro; pero, por otro lado, en términos de bienestar, el resultado es beneficioso para la colectividad, puesto que se reducen los costes y luego los precios (adaptación dinámica).

\subsection{Monopsonio }
¿Qué pasaría se el poder de mercado se ostentase en la otra \textit{cara} de este? Es decir, cómo funcionaría un mercado en el que solo existe un único comprador del bien. Este sería el caso de un monopsonio.

Aunque para una caracterización completa debemos analizar tanto la estructura del mercado en la que ostenta el papel de monpsonista como a la estructura del mercado en la que participa como oferente,\footnote{
Cabe distinguir dos escenarios para el monopsonista en función de la tecnología de producción de la industria competitiva en la que adquiere el factor:
\begin{la}
    \item Rendimientos constantes a escala. En cuyo caso la curva de oferta del factor será infinitamente elástica, pudiendo comprar todas las unidades que desee siempre que pague un precio igual al CMg. En este caso, el gasto medio para el monopsonista es exactamente igual al gasto marginal; y,
    \item Rendimientos decrecientes a escala. En cuyo caso la cuerva de oferta de la industria competitiva tendrá pendiente positiva y la adquisición de una unidad adicional obligará al monopsonista a pagar un mayor precio, tanto por ésta como por las que adquiría con anterioridad.
\end{la}
Por ello, hay que distinguir entre la curva de oferta del factor y la curva del gasto marginal del monopsonio, al igual que ocurría en el caso del monopolio, donde se distinguió entre la curva de demanda y la de ingreso marginal. En este caso, la curva de gasto marginal será siempre superior a su curva de gasto medio, difiriendo en lo que aumenta el precio por todas las unidades afectadas al adquirir el monopsonista una unidad adicional. Asimismo, por lo que se refiere al ingreso marginal del monopsonista, también cabe distinguir dos casos:
\begin{la}
    \item El monopsonista vende el producto en un mercado competitivo. En este caso, carece de poder de mercado respecto a los consumidores y el ingreso marginal coincide con el precio; y,
    \item El monopsonista tiene poder de mercado. Por ello, su ingreso marginal no coincidirá con la curva de demanda de los consumidores.
\end{la}
} supongamos que se trata de un mercado competitivo, en el que se produce un bien intermedio que tiene un único comprador (monopsonio), que lo transforma en un bien que vende en un mercado perfectamente competitivo. Asimismo se supone que los costes de producción del monopsonista están compuestos exclusivamente por los desembolsos monetarios realizados para obtener el bien intermedio. En tal caso, su coste medio o unitario no es más que el precio que debe pagar a sus proveedores. Por lo tanto, este bien intermedio se utiliza como único factor de producción y la decisión del monopsonista relativa al número óptimo de unidades a comprar se determina siguiendo el mismo principio que el monopolio: irá adquiriendo unidades siempre que el valor de su ingreso marginal, determinado por la demanda de consumidores, sea superior al coste marginal de la última unidas adquirida.

\eqblock{
\begin{aligned}
    &\max_{z}\Pi(z)=\underbrace{P_x^d(x^*)\cdot f(z)}_{\text{\scriptsize Ingresos totales}}-\underbrace{P_z^s(z)\cdot z}_{\text{\scriptsize Costes totales}}\\[1em]
    &\text{CPO:}\qquad \frac{\partial}{\partial z}\Pi(z)=\underbrace{P_x^d(x^*)\cdot\frac{\partial}{\partial z}f(z)}_{IMg_z\equiv  \;P_x^d(x^*)\cdot {PMg}_z}-\underbrace{P_z^s(z)-z\frac{\partial P_z^s(z)}{\partial z}}_{\text{\scriptsize Coste marginal del factor}}=0\\[0.5em]
    &\underset{\text{\scriptsize rescribiendo}}{\Longrightarrow}{CMg}_z=P_z^s(z)+\frac{P_z^s(z)}{P_z^s(z)}\cdot z\frac{\partial P_z^s(z)}{\partial z }\Rightarrow\;P_z^s(z)+\frac{P_z^s(z)}{\varepsilon_{S,p_z}}=\boxed{P_z^s(z)\left(1+\frac{1}{\varepsilon_{S,p_z}}\right)}\\[1em]
    &\underset{\text{\scriptsize En equilibrio}}{\Longrightarrow}\qquad \boxed{{IMg}_z={CMg}_z}\quad \text{y}\quad\underset{\text{\scriptsize Margen sobre el óptimo}}{\boxed{\frac{{CMg}_z-P_z^s}{P_z^s}=\frac{1}{\varepsilon_{S,p_z}}}}
\end{aligned}
}{Problema de maximización del monopsonista. Extensiones: Teoría del Monopolio}

El monopsonista maximiza beneficios eligiendo la cantidad del factor $z$, teniendo en cuenta que al aumentar su demanda eleva el precio de todas las unidades adquiridas. Por ello, el coste marginal del factor excede su precio, y la diferencia depende inversamente de la elasticidad de oferta del factor: cuanto más inelástica es la oferta, mayor es el poder de monopsonio.

Como resultado, el monopsonista no iguala el valor de la productividad marginal del factor a su precio, sino a su coste marginal social, que incorpora el efecto de su propia demanda sobre el precio del input. En consecuencia, el monopsonista adquiere el factor hasta que el ingreso marginal social del factor se iguala a su coste marginal social, lo que implica subutilización del factor respecto al óptimo competitivo.

\subsection{Monopolio Bilateral}
Si ponemos todo lo expuesto hasta ahora en común podemos estudiar brevemente lo que sucede cuando se encuentran en el mercado un comprador monopsonista y un vendedor monopolista. Por los motivos teóricos que se expondrán a continuación, se trata de una estructura de mercado que rara vez se encuentra en la práctica, con la excepción de los mercados de trabajo. Supóngase que:
\begin{la}
    \item \textbf{Monopsonista}. La curva de ingresos marginales del producto para el monopsonista es $R$ y, dado que compraría cantidades en esta curva para precios fijos, $R$ constituye la curva de ingresos medios para el monopolista y, en consecuencia, su derivada $R'$, la curva de ingreso marginal para el monopolista. El monopsonista maximizará sus beneficios en la proyección sobre la curva de costes marginales del monopolista ($C$) de la intersección de su propia curva de costes marginales ($C'$) y su propia curva de ingresos marginales ($R$) ;
    \item \textbf{Monopolista}. Un monopolista tiene la curva de costes marginales $C$. Por lo tanto, con precios fijos, ofertará las cantidades indicadas por esta curva, por lo que la misma puede denominarse como la curva de costes medios para el comprador y, en consecuencia, $C'$ es el coste marginal del bien para el comprador. En este caso, el monopolista maximizará sus beneficios en la proyección sobre su curva de ingresos medios ($R$) de la interesección entre su curva de costes marginales ($C$) y su curva de ingresos marginales ($R'$) haga en su curva de ingresos medios.
\end{la}

\imagenfit{Monopolio_Bilateral.png}{Equilibrio indeterminado en un mercado en monopolio bilateral}{Apuntes ICEX-CECO. Tema 3.A.17}

Como los objetivos son incompatibles (salvo por puro azar), el precio bajo un monopolio bilateral queda indeterminado en el sentido que las condiciones de costes y demanda no son suficientes para determinar el precio y la cantidad.

Una posible vía de solución es combinar los beneficios conjuntos de ambas empresas si no intentasen aprovecharse la una de la otra (i.e., si se contentasen con aprovecharse de sus proveedores y compradores). Si $R$ es la curva de ingreso marginal para el comprador y $C$ es la curva de coste marginal para el vendedor, los beneficios totales conjuntos de las dos empresas serían mayores en el output $0F$ que en cualquier otro. Dado que los beneficios de las dos empresas serían mayores si pudiesen operar donde el coste marginal se iguala con el ingreso marginal, existe un fuerte incentivo para ambas de fusionarse. La dificultad a la hora de identificar ejemplos interesantes de monopolio bilateral se deriva de que constituye una forma inestable de organización.


 












\clearpage
\section{Conclusión} 

\subsection{Extensiones y relación con otras partes del Temario}


\subsection{Opinión}



\subsubsection{Idea final (Salida o cierra)}


\clearpage
\pagenumbering{gobble}
\MyBibliography
\MyBibItem{Autor}{Título}{Año}{DOI -- LINK}




%ANEXOS
\clearpage
\anexo{\textbf{Preguntas Test}}
%\input{\TemaCode/questions.tex} 




\end{document}